\documentclass[11pt,a4paper]{article}
\usepackage{fontspec}
\defaultfontfeatures{Ligatures=TeX}
\setmainfont{Times New Roman}
\usepackage[main=russian,english]{babel}
\usepackage[margin=2cm]{geometry}
\usepackage[most]{tcolorbox}
\usepackage{enumitem}
\usepackage{hyperref}
\usepackage{microtype}
\hypersetup{colorlinks=true,linkcolor=blue,urlcolor=blue,pdftitle={Сценарий заседания: Введение в медицинскую статистику}}
\definecolor{headingblue}{HTML}{013362}
\tcbset{
  colback=gray!5, colframe=headingblue!60!black, boxrule=0.4pt,
  arc=2pt, left=6pt, right=6pt, top=4pt, bottom=4pt,
  toptitle=3pt, bottomtitle=3pt, fonttitle=\bfseries,
  breakable
}
\newtcolorbox{onslide}{title={На слайде}}
\newtcolorbox{jamovi}{title={Демонстрация в Jamovi},colback=blue!4,colframe=blue!50!black}
\newtcolorbox{explain}{title={Пояснение},colback=yellow!5,colframe=yellow!60!black}
\setlength{\parindent}{0pt}
\sloppy
\hfuzz=32pt
\setlength{\parskip}{4pt}
\newcommand{\JamoviURL}{https://www.jamovi.org/download.html}
\newcommand{\DataURL}{https://archive.ics.uci.edu/dataset/45/heart+disease}
\newcommand{\RepoURL}{https://github.com/ci-c/SNK-jamovi}

\begin{document}

\begin{center}
{\Large\bfseries Сценарий заседания СНК\\
«Введение в медицинскую статистику»}
\end{center}


\begin{tcolorbox}[title={Ссылки для скачивания}]
Программа Jamovi: \url{\JamoviURL}. Набор данных: \url{\DataURL}.
Репозиторий заседания: \url{\RepoURL}.
\end{tcolorbox}

\begin{tcolorbox}[title={Переменные, с которыми работаем}]
Непрерывные переменные: \texttt{age}, \texttt{trestbps}, \texttt{chol}, \texttt{thalach} и \texttt{oldpeak}. Категориальные переменные: \texttt{sex}, \texttt{cp}, \texttt{fbs}, \texttt{restecg}, \texttt{exang}, \texttt{slope}, \texttt{ca}, \texttt{thal} и \texttt{target}. В расчётах участвуют \texttt{age}, \texttt{chol}, \texttt{thalach} и \texttt{oldpeak}; в сравнениях групп и таблицах сопряжённости участвуют \texttt{sex}, \texttt{cp} и \texttt{target}; в регрессии дополнительно используется \texttt{sex}.
\end{tcolorbox}



\begin{tcolorbox}[title={Карта критериев в Jamovi}]
\begin{itemize}\tiny
\item \textbf{Описательные показатели}. Путь: Analyses $\to$ Exploration $\to$ Descriptives. Документация: \texttt{jg\_11\_descriptive-analyses.html}.
\item \textbf{Шапиро--Уилк}. Путь: Descriptives $\to$ Statistics $\to$ Shapiro-Wilk. Документация: \texttt{jg\_11\_descriptive-analyses.html}.
\item \textbf{Независимый $t$-тест}. Путь: Analyses $\to$ T-Tests $\to$ Independent Samples T-Test. Документация: \texttt{jg\_21\_t-test-independent.html}.
\item \textbf{Парный $t$-тест}. Путь: Analyses $\to$ T-Tests $\to$ Paired Samples T-Test. Документация: \texttt{jg\_22\_t-test-paired.html}.
\item \textbf{Одновыборочный $t$-тест}. Путь: Analyses $\to$ T-Tests $\to$ One Sample T-Test. Документация: \texttt{jg\_23\_t-test-onesample.html}.
\item \textbf{Манн--Уитни}. Опция Mann-Whitney U в Independent Samples T-Test. Документация: \texttt{s2j\_nonParametric.html}.
\item \textbf{Уилкоксон}. Опция Wilcoxon в Paired Samples T-Test. Документация: \texttt{s2j\_nonParametric2.html}.
\item \textbf{Однофакторный ANOVA}. Путь: Analyses $\to$ ANOVA $\to$ One-Way ANOVA. Документация: \texttt{jg\_31\_unianova.html}.
\item \textbf{Корреляционная матрица}. Путь: Analyses $\to$ Regression $\to$ Correlation Matrix. Документация: \texttt{jg\_overview.html}.
\item \textbf{Линейная регрессия}. Путь: Analyses $\to$ Regression $\to$ Linear Regression. Документация: \texttt{jg\_42\_regression-linear.html}.
\item \textbf{Логистическая регрессия}. Путь: Analyses $\to$ Regression $\to$ Logistic Regression. Документация: \texttt{jg\_43\_regression-logistic.html}.
\item \textbf{Хи-квадрат}. Путь: Analyses $\to$ Frequencies $\to$ Independent Samples, $\chi^2$ test. Документация: \texttt{s2j\_frequencies.html}.
\item \textbf{Точный критерий Фишера}. Опция Fisher's exact test в Statistics таблицы сопряжённости. Документация: \texttt{s2j\_frequencies.html}.
\item \textbf{Байесовский фактор}. Опция Bayes factor в опциях $t$-тестов. Документация: \texttt{jg\_overview.html}.
\item \textbf{Вычисленные переменные}. Путь: Data $\to$ Compute. Документация: \texttt{data\_2\_computed\_variables.html}.
\item \textbf{Фильтры}. Путь: Data $\to$ Filter. Документация: \texttt{data\_6\_filtering\_data.html}.
\end{itemize}
Полные адреса документов расположены по префиксу \url{https://docs.jamovi.org/}.
\end{tcolorbox}

\section*{Структура заседания}

\begin{itemize}
\item Вступление, слайды 1--8: титул, материалы, план, Jamovi и ссылки.
  \begin{itemize}
    \item Титул и материалы, слайды 1--2.
    \item План заседания, слайд 3.
    \item Цель статистики и подготовка к работе, слайды 4--8.
  \end{itemize}
\item Данные, слайды 9--14.
  \begin{itemize}
    \item Импорт и структура набора, слайды 10--11.
    \item Подготовка данных, пропуски и словарь данных, слайды 12--13.
    \item Типы переменных, слайд 14.
  \end{itemize}
\item Описательная статистика и нормальность, слайды 15--22.
  \begin{itemize}
    \item Показатели и правила записи, слайды 16--18.
    \item Асимметрия, эксцесс и преобразования, слайд 19.
    \item Критерий Шапиро--Уилка и выбор критерия, слайды 21--22.
  \end{itemize}
\item Сравнение групп, слайды 23--31.
  \begin{itemize}
    \item Выборки и $t$-критерий, слайды 24--26.
    \item $p$-значение и доверительный интервал, слайды 27--28.
    \item Памятка, эффект, мощность и дополнительные критерии, слайды 29--31.
  \end{itemize}
\item Корреляция и регрессия, слайды 32--38.
  \begin{itemize}
    \item Корреляция, слайды 33--34.
    \item Линейная и множественная регрессия, проверки модели, слайды 35--38.
  \end{itemize}
\item Категориальные данные, слайды 39--42.
  \begin{itemize}
    \item Таблица сопряжённости, хи-квадрат, Фишер и логистическая регрессия, слайды 40--42.
  \end{itemize}
\item Графики, слайды 43--51.
  \begin{itemize}
    \item Базовые графики и графики из статей, слайды 44--50.
    \item Подписи графиков, слайд 51.
  \end{itemize}
\item Байесовский статистический вывод, слайды 52--66.
  \begin{itemize}
    \item Разные вопросы и формула Байеса, слайды 53--54.
    \item Байесовский фактор и оценивание, слайды 55--58.
    \item Шансы, чувствительность, MCMC и ограничения, слайды 59--65.
  \end{itemize}
\item Чтение клинической статьи, слайды 67--88.
  \begin{itemize}
    \item Дизайн, поток участников и таблицы, слайды 67--70.
    \item Графики, ограничения и дизайн исследования, слайды 71--76.
    \item ITT, статистический план, эффекты и применимость, слайды 77--84.
    \item Чек-лист и ответы, слайды 85--87.
  \end{itemize}
\item ИИ, слайды 89--91.
\item Итоги, слайды 92--95.
\end{itemize}

\begin{tcolorbox}[title={Условные обозначения в формулах}]
\begin{itemize}\footnotesize
\item $x_i$ — отдельное наблюдение, $\bar{x}$ — выборочное среднее, $n$ — объём выборки, $s$ — стандартное отклонение.
\item $Q_1$ и $Q_3$ — первый и третий квартили, $IQR=Q_3-Q_1$ — межквартильный размах.
\item $\mu_0$ — проверяемое значение среднего, $\alpha$ — уровень значимости, $p$ — $p$-значение.
\item $H_0$ и $H_1$ — нулевая и альтернативная гипотезы, $D$ — данные, $T$ — статистика критерия, $t_{\text{набл}}$ — наблюдаемое значение статистики.
\item $\chi^2$ — статистика критерия хи-квадрат, $df$ — число степеней свободы, $O$ и $E$ — наблюдаемая и ожидаемая частоты, $V$ — коэффициент Крамера.
\item $r$ — коэффициент корреляции Пирсона, $\rho$ — коэффициент корреляции Спирмена, $z$ — значение после преобразования Фишера.
\item $B$ и $\beta$ — коэффициенты регрессии, $Y$ — отклик, $X$ — предиктор, $\varepsilon$ — случайная ошибка, $R^2$ — коэффициент детерминации.
\item $d$ — стандартизованная разность средних Коэна, $g$ — стандартизованная разность Хеджеса, $J$ — поправка на объём выборки, $\beta$ — вероятность ошибки второго рода.
\item $\mathrm{BF}_{10}$ — байесовский фактор в пользу $H_1$, $\theta$ — параметр модели, $p(D\mid H)$ — маргинальное правдоподобие.
\item $\mathrm{PPV}$ — прогностическая ценность положительного результата, $\mathrm{OR}$ — отношение шансов, $\mathrm{HR}$ — отношение рисков, $\mathrm{ARR}$ — абсолютное снижение риска, $\mathrm{NNT}$ — число пациентов, которых нужно пролечить.
\end{itemize}
\end{tcolorbox}

\section*{Как читать сценарий}

Для каждого слайда приведены отдельные блоки. Блок «На слайде» описывает содержимое,
которое видит аудитория. Блок «Демонстрация в Jamovi» описывает действия на компьютере
ведущего; в нём указаны только клавиатурные команды: \texttt{Tab}, \texttt{Shift+Tab},
\texttt{Enter}, \texttt{Space}, стрелки и \texttt{Alt} с последующей буквой. Блок
«Пояснение» содержит нейтральные формулировки для устного выступления и указания на
то, на что именно смотреть на изображении. Задания для аудитории в материале
отсутствуют.

\section*{Сценарий по слайдам}


\subsection*{Слайд 1. Введение в медицинскую статистику}
\begin{onslide}
На экране находится титульный слайд с названием заседания, именем ведущего и названием кафедры.
\end{onslide}
На слайде показан титульный лист с названием заседания, именем ведущего и названием кафедры. Ведущий зачитывает название по мере необходимости и сохраняет его без изменений. Титульный слайд задаёт тему и организационную рамку для всего изложения. После титула аудитория переходит к материалам, которые понадобятся на протяжении всего заседания.

\subsection*{Слайд 2. Материалы заседания}
\begin{onslide}
На слайде расположен QR-код репозитория. Рядом указано, что презентация, сценарий и датасет хранятся в репозитории заседания.
\end{onslide}
На слайде изображён QR-код репозитория, рядом указано, что презентация, сценарий и датасет хранятся в одном месте. Ведущий сообщает, что ссылка продублирована в посте, и выдерживает паузу 20--30 секунд на сканирование. Единое хранилище важно, чтобы слушатели работали с теми же файлами, что и ведущий. Следующий слайд раскрывает план изложения.

\subsection*{Слайд 3. План заседания}
\begin{onslide}
Слайд содержит план изложения по разделам и диапазонам слайдов: данные и описание (9--22), сравнение групп и связи (23--51), байесовский вывод (52--65), чтение статьи и итоги (66--94). В блоке «Формат» указано, что формально объявлено одно заседание на 60 минут, материал подготовлен с запасом, редактура и урезание будут выполнены позже.
\end{onslide}
На слайде приведён план изложения по темам и блок «Формат». Ведущий сообщает, что формально объявлено одно заседание на 60 минут, материал подготовлен с запасом, а редактура и урезание будут выполнены позже. План охватывает данные, описание, проверку распределения, сравнение групп, связи, графики, байесовский статистический вывод, чтение статьи и ограничения применения ИИ. Ведущий демонстрирует действия, а слушатели повторяют их на открытом датасете UCI. План подготавливает переход к разделу «Введение».

\subsection*{Слайд 4. Введение}
\begin{onslide}
Слайд-разделитель с названием раздела.
\end{onslide}
На слайде показан разделитель раздела «Введение». Ведущий сообщает, что раздел определяет цель заседания и место Jamovi в практической части. Разделитель отделяет организационные слайды от содержательных. Следующий слайд формулирует цель статистики.

\subsection*{Слайд 5. Статистика служит языком научных результатов}
\begin{onslide}
На слайде указано, что статистика выступает языком, на котором наука сообщает результаты. Во втором пункте говорится о рассмотрении инструментов получения таблицы и её чтения. Блок «К концу заседания» сообщает, что слушатель открывает Jamovi, рассчитывает показатели и тесты, читает $p$-значение и доверительный интервал.
\end{onslide}
На слайде указано, что статистика выступает языком, на котором наука сообщает результаты, и что рассматриваются инструменты получения таблицы и её чтения. Блок «К концу заседания» сообщает, что слушатель открывает Jamovi, рассчитывает показатели и тесты, читает $p$-значение и доверительный интервал. Ведущий поясняет, что Jamovi предоставляет готовую информацию, а вывод формул в заседании не планируется. Результатом заседания становится умение получить таблицу и интерпретировать $p$-значение, доверительный интервал и результат байесовского анализа. Следующий слайд знакомит с программой, которая обеспечивает практическую часть.

\subsection*{Слайд 6. Возможности Jamovi и установка программы}
\begin{onslide}
Слайд содержит три блока. Блок «Что это» сообщает, что программа распространяется бесплатно, её исходный код открыт, работа выполняется через меню, знание языка R не требуется, программа формирует таблицы для статьи. Блок «Где скачать» указывает сайт \texttt{jamovi.org} и поддержку Windows, macOS и Linux через flatpak. Блок «Интерфейс» сообщает, что данные расположены слева, результаты справа, меню сверху.
\end{onslide}
На слайде приведены три блока: «Что это», «Где скачать» и «Интерфейс». Ведущий сообщает, что Jamovi распространяется бесплатно, имеет открытый исходный код, работает через меню, не требует знания языка R и формирует таблицы для статьи. Сайт jamovi.org поддерживает Windows, macOS и Linux через flatpak. В интерфейсе данные расположены слева, результаты справа, а меню сверху. Следующий слайд показывает страницу загрузки программы.

\subsection*{Слайд 7. Страница загрузки Jamovi содержит ссылку Download Desktop}
\begin{onslide}
Полноэкранный скриншот страницы загрузки Jamovi.
\end{onslide}
На слайде представлен полноэкранный скриншот страницы загрузки Jamovi. Ведущий демонстрирует снимок и указывает на ссылку Download Desktop как на начало установки. Адрес страницы приведён на следующем слайде. Скриншот позволяет слушателям найти нужную ссылку без дополнительного поиска.

\subsection*{Слайд 8. Ссылки для скачивания программы и данных}
\begin{onslide}
На слайде расположены два QR-кода: для страницы загрузки Jamovi \url{https://www.jamovi.org/download.html} и для набора данных UCI Heart Disease \url{https://archive.ics.uci.edu/dataset/45/heart+disease}.
\end{onslide}
{\small На слайде расположены два QR-кода для страницы загрузки Jamovi \url{https://www.jamovi.org/download.html} и набора данных UCI Heart Disease \url{https://archive.ics.uci.edu/dataset/45/heart+disease}.\
Ведущий сообщает, что оба QR-кода ведут на официальные страницы. Репозиторий заседания указан на слайде 2. Ссылки завершают подготовительный блок, после чего изложение переходит к данным.}

\subsection*{Слайд 9. Данные}
\begin{onslide}
Слайд-разделитель с названием раздела.
\end{onslide}
На слайде показан разделитель раздела «Данные». Ведущий сообщает, что все расчёты выполняются на одном открытом наборе данных. Сначала проверяются структура набора и типы переменных. Разделитель подготавливает импорт файла на следующем слайде.

\subsection*{Слайд 10. После импорта в Jamovi отображаются 303 пациента и 14 колонок}
\begin{onslide}
Полноэкранный скриншот окна Jamovi после загрузки файла \texttt{heart-disease.csv}. Внизу окна видны счётчики наблюдений и переменных.
\end{onslide}
\begin{jamovi}
Выполнить. 1. Запустить Jamovi. 2. Нажать \texttt{Alt}, затем нажать \texttt{F}, затем \texttt{O}; в диалоге указать путь к \texttt{heart-disease.csv} и подтвердить \texttt{Enter}. 3. Показать счётчик внизу окна: 303 наблюдения и 14 переменных.
\end{jamovi}
На слайде показан скриншот окна Jamovi после загрузки файла \texttt{heart-disease.csv}; внизу окна видны счётчики наблюдений и переменных. Ведущий запускает Jamovi, открывает файл через меню и показывает счётчик: 303 наблюдения и 14 переменных. Ведущий проверяет, что импортированы 303 строки и 14 колонок. Типичная ошибка заключается в открытии CSV двойным кликом в табличном редакторе. Файл открывается через меню \texttt{Open} в Jamovi, после чего можно переходить к переменным.

\subsection*{Слайд 11. Какие переменные содержатся в наборе данных}
\begin{onslide}
На слайде перечислены непрерывные переменные \texttt{age}, \texttt{chol}, \texttt{thalach} и \texttt{oldpeak}; категориальные переменные \texttt{cp} с четырьмя типами боли, \texttt{sex}, \texttt{fbs}, \texttt{exang}; исход \texttt{target} принимает значение 1 при наличии болезни и 0 при её отсутствии. Блок «Открытые данные» сообщает, что набор UCI Heart Disease содержит 303 пациента.
\end{onslide}
На слайде перечислены непрерывные переменные \texttt{age}, \texttt{chol}, \texttt{thalach} и \texttt{oldpeak}, а также категориальные переменные \texttt{cp} с четырьмя типами боли, \texttt{sex}, \texttt{fbs} и \texttt{exang}. Исход \texttt{target} принимает значение 1 при наличии болезни и 0 при её отсутствии, а блок «Открытые данные» сообщает, что набор содержит 303 пациента. Ведущий дополняет список категориальных переменных: \texttt{restecg}, \texttt{slope}, \texttt{ca} и \texttt{thal}, а также \texttt{target}. В файле имеются шесть пустых ячеек, которые Jamovi импортирует как пропуски. Следующий слайд показывает, как задать типы переменных.

\subsection*{Слайд 12. Подготовка данных в Jamovi}
\begin{onslide}
На слайде описаны фильтры, перекодирование и вычисленные переменные. Документация Jamovi описывает measure type как подсказку. В большинстве анализов программа приводит переменную к нужному типу, а явная установка типа упрощает проверку данных.
\end{onslide}
\begin{jamovi}
Выполнить. 1. Нажать \texttt{Alt}, затем нажать \texttt{D}; \texttt{Tab} к разделу \texttt{Filter}, \texttt{Enter} задаёт условие отбора. 2. В редакторе переменных выбрать колонку и через \texttt{Recode} создать новую категорию. 3. Через \texttt{Compute} создать переменную по формуле, например $\texttt{chol\_mmol}=\texttt{chol}/38{,}67$.\par Формула вычисленной переменной показывает преобразование исходной колонки. Смысл операции состоит в получении новой величины с явно заданным правилом.
\end{jamovi}
Фильтр не удаляет исходные строки, поэтому результат можно проверить. Перекодирование полезно, когда исходные коды неудобны для чтения таблицы. Вычисленная переменная позволяет перейти к другой единице измерения и сохранить формулу в сценарии.

\subsection*{Слайд 13. Пропуски и словарь данных}
\begin{onslide}
На слайде описаны три механизма пропусков. MCAR означает, что вероятность пропуска не зависит от значений. MAR означает, что вероятность пропуска зависит от наблюдаемых переменных. MNAR означает, что вероятность пропуска зависит от ненаблюдаемых значений. Описаны полные случаи, одиночное вменение и множественное вменение с оговорками о смещении и неопределённости. Словарь данных содержит название колонки, тип, единицы измерения, допустимые значения и правило кодировки пропусков.
\end{onslide}
Ведущий объясняет, что MCAR, MAR и MNAR являются допущениями о механизме пропусков. Выбор способа зависит от механизма, объёма неполноты и модели. Исключение неполных наблюдений может быть неэффективным и смещённым. Простая замена средним или медианой не сохраняет неопределённость. Множественное вменение требует правдоподобной модели и анализа чувствительности. Словарь данных делает анализ воспроизводимым, потому что фиксирует смысл каждой колонки.

\subsection*{Слайд 14. Типы переменных задаются до выполнения тестов}
\begin{onslide}
Слева приведены списки: номинальные переменные \texttt{sex}, \texttt{cp}, \texttt{fbs}, \texttt{restecg}, \texttt{exang}, \texttt{slope}, \texttt{ca}, \texttt{thal}, \texttt{target}; непрерывные переменные \texttt{age}, \texttt{trestbps}, \texttt{chol}, \texttt{thalach}, \texttt{oldpeak}. Справа показан фрагмент вкладки Variables. Внизу описана установка \texttt{Measure type = Nominal} и её последствия.
\end{onslide}
\begin{jamovi}
Выполнить. 1. Нажать \texttt{Alt}, затем нажать \texttt{V}; откроется вкладка \texttt{Variables}. 2. Выделить строку переменной и нажать \texttt{Enter}. 3. \texttt{Tab} к полю \texttt{Measure type}, \texttt{Space} раскрывает список, стрелкой выбрать \texttt{Nominal}, \texttt{Enter} подтверждает. 4. Повторить для всех переменных из списка номинальных.
\end{jamovi}
На слайде слева приведены списки номинальных и непрерывных переменных, а справа показан фрагмент вкладки Variables. Номинальные переменные включают \texttt{sex}, \texttt{cp}, \texttt{fbs}, \texttt{restecg}, \texttt{exang}, \texttt{slope}, \texttt{ca}, \texttt{thal} и \texttt{target}; при этом \texttt{slope} имеет естественный порядок, а \texttt{ca} является счётной величиной, а непрерывные переменные включают \texttt{age}, \texttt{trestbps}, \texttt{chol}, \texttt{thalach} и \texttt{oldpeak}. В Jamovi выделяется строка переменной, после чего для неё устанавливается \texttt{Measure type = Nominal}. Ведущий проверяет значки у колонок: у \texttt{sex} и \texttt{target} категориальный, у \texttt{age} и \texttt{chol} шкала. Дальнейшие расчёты выполняются только после проверки типов, с чего начинается описательная статистика.

\subsection*{Слайд 15. Описательная статистика}
\begin{onslide}
Слайд-разделитель с названием раздела.
\end{onslide}
На слайде показан разделитель раздела «Описательная статистика». Ведущий сообщает, что здесь появляется первая таблица и правила её формулировки. Разделитель отделяет подготовку данных от новых расчётов. Следующий слайд объясняет выбор меры центра и меры разброса.

\subsection*{Слайд 16. Описательная статистика описывает центр и разброс}
\begin{onslide}
В блоке «Центр» указано, что среднее применяется при симметричном распределении, а медиана применяется при скошенном распределении и при выбросах. В блоке «Разброс» указано, что стандартное отклонение указывается со средним, а межквартильный размах указывается с медианой. Ниже сказано, что разброс дополняет среднее: одинаковые средние значения могут соответствовать разным данным.
\end{onslide}
На слайде в блоке «Центр» указано, что среднее применяется при симметричном распределении, а медиана применяется при скошенном распределении и при выбросах. В блоке «Разброс» указано, что стандартное отклонение указывается со средним, а межквартильный размах указывается с медианой. Ведущий поясняет, что разброс дополняет среднее, поскольку одинаковые средние значения могут соответствовать разным данным. Выбор меры центра и меры разброса зависит от вида распределения. Следующий слайд показывает расчёт этих показателей в Jamovi.

\subsection*{Слайд 17. Как рассчитать описательные показатели в Jamovi}
\begin{onslide}
Слайд содержит последовательность: открыть \texttt{Exploration} и выбрать \texttt{Descriptives}; перенести \texttt{age}, \texttt{chol} и \texttt{thalach} в поле \texttt{Variables}; в разделе \texttt{Statistics} включить Mean, SD, Median, Min и Max, для квартилей задать Percentiles 25 и 75. Справа показан фрагмент таблицы результатов.
\end{onslide}
\begin{jamovi}
Выполнить. 1. Нажать \texttt{Alt}, затем нажать \texttt{A}; \texttt{Tab} до \texttt{Exploration}, \texttt{Enter}, стрелкой выбрать \texttt{Descriptives}, \texttt{Enter}. 2. Выделить переменные и перенести их в поле \texttt{Variables} нажатием \texttt{Enter}. 3. \texttt{Tab} к опциям, \texttt{Space} включает нужные показатели.
\par Формулы: $\bar{x}=\frac{1}{n}\sum x_i$, $s=\sqrt{\frac{1}{n-1}\sum(x_i-\bar{x})^2}$, $IQR=Q_3-Q_1$. Смысл: среднее и стандартное отклонение описывают центр и разброс при симметричном распределении, а медиана и межквартильный размах устойчивы к выбросам.
\end{jamovi}
На слайде приведена последовательность расчёта в Jamovi: открыть \texttt{Exploration} и \texttt{Descriptives}, перенести \texttt{age}, \texttt{chol} и \texttt{thalach} в поле \texttt{Variables}. В разделе \texttt{Statistics} включаются Mean, SD, Median, Min и Max, а для квартилей задаются Percentiles 25 и 75. Справа показан фрагмент таблицы результатов. Ведущий использует формулировку «Возраст пациентов составил 54,4 $\pm$ 9,0 года ($n = 303$)». Медиана и межквартильный размах приведены отдельно при необходимости. Объём выборки указывается в каждой формулировке. Следующий слайд разбирает компоненты этой записи.

\subsection*{Слайд 18. Описательный результат записывают одной формулировкой}
\begin{onslide}
Слайд содержит готовую формулировку: «Возраст пациентов составил 54,4 $\pm$ 9,0 года ($n = 303$)». Медиана и межквартильный размах приведены отдельно при необходимости. Ниже указано, что среднее и стандартное отклонение применяются при симметричном распределении, медиана и межквартильный размах применяются при скошенном распределении или при выбросах, объём выборки указывается всегда.
\end{onslide}
На слайде приведена готовая формулировка «Возраст пациентов составил 54,4 $\pm$ 9,0 года ($n = 303$)». Медиана и межквартильный размах приведены отдельно при необходимости. Ниже указано, что среднее и стандартное отклонение применяются при симметричном распределении, а медиана и межквартильный размах применяются при скошенном распределении или при выбросах. Объём выборки указывается всегда. Ведущий разбирает компоненты формулировки и отмечает типичную ошибку использования интервала «среднее $\pm$ SD» при скошенном распределении. Следующий раздел посвящён проверке нормальности.

\subsection*{Слайд 19. Асимметрия, эксцесс и преобразования}
\begin{onslide}
На слайде описаны коэффициенты асимметрии и эксцесса. Стандартная ошибка асимметрии приблизительно равна $\sqrt{6/n}$, стандартная ошибка эксцесса также приблизительно равна $\sqrt{24/n}$. Логарифмирование $y'=\ln(y)$ применяется к положительным значениям при правосторонней асимметрии.
\end{onslide}
\begin{jamovi}
Выполнить. 1. В \texttt{Descriptives} включить \texttt{Skewness} и \texttt{Kurtosis}. 2. Через \texttt{Compute} создать переменную с формулой $\ln(y)$.\par Формула стандартной ошибки асимметрии равна $\sqrt{6/n}$, формула стандартной ошибки эксцесса равна $\sqrt{24/n}$. Смысл преобразования состоит в приближении распределения к нормальному виду.
\end{jamovi}
Коэффициенты асимметрии и эксцесса дополняют гистограмму числовой оценкой формы распределения. Логарифмирование применимо только к положительным значениям. После преобразования вывод относится к логарифмированной шкале, поэтому его нужно явно указывать в отчёте.

\subsection*{Слайд 20. Нормальность}
\begin{onslide}
Слайд-разделитель с названием раздела.
\end{onslide}
На слайде показан разделитель раздела «Нормальность». Ведущий сообщает, что проверка предпосылок помогает выбрать метод с учётом вопроса и дизайна. Разделитель подготавливает переход от описания к проверке гипотез. Следующий слайд демонстрирует критерий Шапиро--Уилка.

\subsection*{Слайд 21. Проверка нормальности в Jamovi с помощью критерия Шапиро–Уилка}
\begin{onslide}
Слайд содержит последовательность: открыть \texttt{Exploration} и \texttt{Descriptives}; перенести \texttt{thalach} и \texttt{oldpeak}; в разделе \texttt{Statistics} включить \texttt{Normality}. Блок «Что читается» сообщает, что критерий возвращает статистику $W$ и $p$-значение; при $p < 0{,}001$ гипотеза о нормальности отвергается, после чего рассматриваются гистограмма и Q-Q-график.
\end{onslide}
\begin{jamovi}
Выполнить. 1. В открытом \texttt{Descriptives} перейти \texttt{Tab} к \texttt{Plots}, \texttt{Enter}; \texttt{Space} включает \texttt{Histogram} и \texttt{Q-Q plot}. 2. \texttt{Tab} к \texttt{Statistics}, \texttt{Enter}; \texttt{Space} включает \texttt{Shapiro-Wilk}.
\par Формула: $W=\frac{(\sum a_i x_{(i)})^2}{\sum(x_i-\bar{x})^2}$. Смысл: статистика сравнивает упорядоченные наблюдения с ожидаемыми при нормальном распределении; малое $p$-значение указывает на отклонение от нормальности.
\end{jamovi}
На слайде приведена последовательность проверки: открыть \texttt{Exploration} и \texttt{Descriptives}, перенести \texttt{thalach} и \texttt{oldpeak} и включить \texttt{Normality} в разделе \texttt{Statistics}. Блок «Что читается» сообщает, что критерий возвращает статистику $W$ и $p$-значение. При $p < 0{,}001$ гипотеза о нормальности отвергается, после чего рассматриваются гистограмма и Q-Q-график. Ведущий сравнивает гистограммы с результатом теста: обе переменные отвергают нормальность при $p < 0{,}05$. На Q-Q-графике точки около линии соответствуют нормальности, а систематические дуги указывают на скошенность. Следующий слайд объясняет, как результат влияет на выбор критерия.

\subsection*{Слайд 22. Выбор критерия зависит от вопроса и предпосылок модели}
\begin{onslide}
На слайде указано, что выбор метода определяется исследовательским вопросом, типом исхода, дизайном и предпосылками конкретной модели. Нормальность оценивают применительно к модели и графическому анализу. Результат Шапиро--Уилка нельзя использовать как единственное правило выбора.
\end{onslide}
На слайде указано, что выбор метода определяется исследовательским вопросом, типом исхода, дизайном и предпосылками модели. Независимый $t$-тест требует независимых наблюдений, оценки выбросов и распределения исхода внутри групп. Парный $t$-тест оценивает распределение парных разностей. Корреляция Пирсона предполагает линейность и чувствительна к выбросам. Регрессия оценивает линейность, гомоскедастичность и структуру остатков. Ведущий обсуждает альтернативные критерии и переходит к сравнению групп.

\subsection*{Слайд 23. Сравнение групп}
\begin{onslide}
Слайд-разделитель с названием раздела.
\end{onslide}
На слайде показан разделитель раздела «Сравнение групп». Ведущий сообщает, что первый вопрос состоит в том, какие выборки рассматриваются, независимые или связанные. Разделитель подготавливает выбор критерия для сравнения. Следующий слайд различает связанные и несвязанные выборки.

\subsection*{Слайд 24. Выборки бывают связанными и несвязанными}
\begin{onslide}
В левом блоке указано, что разные пациенты образуют несвязанные выборки, например группы по полу, и используется \texttt{Independent Samples T-Test}. В правом блоке указано, что одни и те же пациенты с измерениями до и после образуют связанные выборки и используется \texttt{Paired Samples T-Test}. Ниже перечислены три вопроса: число групп, независимость групп и нормальность распределения.
\end{onslide}
На слайде в левом блоке указано, что разные пациенты образуют несвязанные выборки, например группы по полу, и используется \texttt{Independent Samples T-Test}. В правом блоке указано, что одни и те же пациенты с измерениями до и после образуют связанные выборки и используется \texttt{Paired Samples T-Test}. Ниже перечислены три вопроса: число групп, независимость групп и нормальность распределения. Ведущий сообщает, что в данном наборе данных парных измерений нет. Пример связанных выборок демонстрируется на синтетических данных, а следующий слайд показывает результат сравнения несвязанных групп.

\subsection*{Слайд 25. Independent Samples T-Test: сравнение chol по sex}
\begin{onslide}
Полноэкранный скриншот результата \texttt{Independent Samples T-Test} для переменной \texttt{chol} с группировкой по \texttt{sex}.
\end{onslide}
\begin{jamovi}
Выполнить. 1. Нажать \texttt{Alt}, затем нажать \texttt{A}; \texttt{Tab} до \texttt{T-Tests}, \texttt{Enter}; выбрать \texttt{Independent Samples T-Test}, \texttt{Enter}. 2. \texttt{Tab} к Dependent Variable, выбрать \texttt{chol}, \texttt{Enter}; к Grouping Variable, выбрать \texttt{sex}, \texttt{Enter}. 3. \texttt{Tab} к опциям, \texttt{Space} включает \texttt{Descriptives}, \texttt{Normality test}, \texttt{Homogeneity test} и \texttt{Mean difference}.
\par Формула: $t=\frac{\bar{x}_1-\bar{x}_2}{s_p\sqrt{1/n_1+1/n_2}}$. Смысл: отношение разности средних к её стандартной ошибке; $p$-значение показывает вероятность такой или более крайней разности при нулевой гипотезе.
\end{jamovi}
На слайде показан полноэкранный скриншот результата \texttt{Independent Samples T-Test} для переменной \texttt{chol} с группировкой по \texttt{sex}. Ведущий открывает \texttt{T-Tests}, выбирает \texttt{Independent Samples T-Test}, задаёт \texttt{chol} в поле Dependent Variable и \texttt{sex} в поле Grouping Variable. В опциях включаются \texttt{Descriptives}, \texttt{Normality test}, \texttt{Homogeneity test} и \texttt{Mean difference}. На скриншоте показана строка с $t = 3{,}54$, $df = 301$, $p < 0{,}001$ и разницей средних 22,2. Следующий слайд объясняет, как читать эти числа.

\subsection*{Слайд 26. Как читать результат t-критерия}
\begin{onslide}
На слайде разобраны согласованные результаты двух вариантов теста. Классический $t$-тест с объединённой дисперсией даёт $t = 3{,}54$, $df = 301$, $p < 0{,}001$ и 95\% доверительный интервал [9,8; 34,5] мг/дл. Критерий Левена даёт $p < 0{,}001$, поэтому основным считается тест Уэлча: разность средних 22,15 мг/дл, $t = 3{,}06$, $df \approx 136{,}37$, $p = 0{,}00263$, 95\% доверительный интервал [7,9; 36,4] мг/дл.
\end{onslide}
На слайде разобраны классический и уточнённый результаты. Классический $t$-тест даёт $t = 3{,}54$, $df = 301$, $p < 0{,}001$ и доверительный интервал [9,8; 34,5] мг/дл. Из-за неоднородности дисперсий основным считается тест Уэлча: разность 22,15 мг/дл, $t = 3{,}06$, $df \approx 136{,}37$, $p = 0{,}00263$, доверительный интервал [7,9; 36,4] мг/дл. Ведущий применяет правило: если доверительный интервал разности не пересекает ноль, различие значимо. Следующий слайд уточняет смысл $p$-значения.

\subsection*{Слайд 27. Что показывает p-значение}
\begin{onslide}
В блоке «Определение» сказано, что $p$-значение — вероятность при условии истинности нулевой гипотезы и заданной статистической модели получить наблюдаемое значение статистики или более экстремальное значение. Блок «Смысл p-значения» сообщает, что $p$-значение описывает данные при нулевой гипотезе и модели, что $p$ не равно вероятности истинности $H_0$ и не измеряет величину эффекта. Блок «Свойства p-значения» сообщает, что $p$-значение используют для решения об отвержении нулевой гипотезы и что оно зависит от объёма выборки.
\end{onslide}
На слайде в блоке «Определение» сказано, что $p$-значение — вероятность при условии истинности нулевой гипотезы и заданной модели получить наблюдаемое значение статистики или более экстремальное значение. Блок «Смысл p-значения» сообщает, что $p$-значение описывает данные при нулевой гипотезе и модели, что $p$ не равно вероятности истинности $H_0$ и не измеряет величину эффекта. Блок «Свойства p-значения» сообщает, что $p$-значение используют для решения об отвержении нулевой гипотезы и что оно зависит от объёма выборки. Ведущий подчёркивает, что величина $p$ описывает данные при нулевой гипотезе. Следующий слайд показывает, как дополнительно к $p$ читать доверительный интервал.

\subsection*{Слайд 28. Доверительный интервал показывает согласующиеся значения}
\begin{onslide}
Блок «Что это» сообщает, что 95\% доверительный интервал получен процедурой, которая при многократном повторении исследования накрывает истинное значение параметра примерно в 95\% случаев при выполнении предпосылок. Ниже указано, что для того же критерия интервал разности, не пересекающий ноль, соответствует значимому различию, и что более узкий интервал соответствует более точной оценке.
\end{onslide}
На слайде в блоке «Что это» сообщается, что доверительный интервал образует диапазон значений, согласующихся с данными при 95\% доверительном уровне. Для того же критерия интервал разности, не пересекающий ноль, соответствует значимому различию. Более узкий интервал соответствует более точной оценке. Ведущий указывает, что тот же интервал используется в таблицах $t$-критерия, регрессии и байесовского оценивания с разным смыслом. Памятка на следующем слайде помогает выбрать критерий для каждой задачи.

\subsection*{Слайд 29. Памятка по выбору критерия}
\begin{onslide}
Таблица связывает задачу и критерий. Для сравнения двух групп с непрерывным исходом выбирается \texttt{Independent T-Test} с вариантом Уэлча. Для рангового исхода выбирается \texttt{Mann--Whitney U}. Для парных непрерывных разностей выбирается \texttt{Paired T-Test}, а для ранговых парных разностей критерий Уилкоксона. Для сравнения трёх и более групп выбирается \texttt{One-Way ANOVA}. Для связи двух переменных выбирается корреляция Пирсона или Спирмена. Для фактора с одной непрерывной переменной выбирается $t$-тест или ANOVA. Для ассоциации категорий выбирается \texttt{Chi-square test}. Для прогноза непрерывной зависимости выбирается \texttt{Linear regression}. Для прогноза бинарной зависимости выбирается \texttt{Logistic regression}.
\end{onslide}
На слайде приведена таблица, которая связывает задачу и критерий. Для сравнения двух групп с непрерывным исходом выбирается \texttt{Independent T-Test} с вариантом Уэлча, а для рангового исхода выбирается \texttt{Mann--Whitney U}. Для парных непрерывных разностей выбирается \texttt{Paired T-Test}, а для ранговых разностей критерий Уилкоксона, а для сравнения трёх и более групп выбирается \texttt{One-Way ANOVA}. Для связи двух переменных выбирается корреляция Пирсона или Спирмена, для фактора с одной непрерывной переменной выбирается $t$-тест или ANOVA, для ассоциации категорий выбирается \texttt{Chi-square test}, для прогноза непрерывной зависимости выбирается \texttt{Linear regression}, а для прогноза бинарной зависимости выбирается \texttt{Logistic regression}. Таблица является раздаточным материалом, и ведущий отдельно указывает различие между статистической и клинической значимостью: $p < 0{,}001$ на 303 пациентах не определяет клиническую важность разницы. Следующий раздел переходит к связям между переменными.

\subsection*{Слайд 30. Величина эффекта, мощность и объём выборки}
\begin{onslide}
На слайде приведены формулы $d$ Коэна и $g$ Хеджеса. Мощность равна $1-\beta$ и представляет вероятность отвергнуть нулевую гипотезу при заданном эффекте. Объём выборки рассчитывается по ожидаемому эффекту, мощности и уровню значимости.
\end{onslide}
\begin{jamovi}
Выполнить. В таблице \texttt{Independent Samples T-Test} включить \texttt{Effect size}.\par Формула $d$ Коэна равна $(\bar{x}_1-\bar{x}_2)/s$, формула $g$ Хеджеса включает поправку на размер выборки. Смысл величины эффекта состоит в оценке практической значимости различия.
\end{jamovi}
Величина эффекта дополняет $p$-значение и зависит от объёма выборки в меньшей степени. Мощность описывает вероятность обнаружить реальный эффект. Планирование объёма выборки выполняется до сбора данных и должно опираться на клинически осмысленный эффект.

\subsection*{Слайд 31. Дополнительные критерии для средних}
\begin{onslide}
Таблица связывает критерий и задачу. Одновыборочный $t$-тест сравнивает среднее с фиксированным значением. Парный $t$-тест сравнивает разности парных измерений с нулём. Критерий Манна--Уитни сравнивает ранги несвязанных выборок. Критерий Уилкоксона сравнивает ранги парных разностей. ANOVA сравнивает средние трёх и более групп.
\end{onslide}
\begin{jamovi}
Выполнить. 1. Открыть \texttt{T-Tests} через \texttt{Alt} и \texttt{A}; выбрать \texttt{One Sample T-Test} или \texttt{Paired Samples T-Test}. 2. Для несвязанных выборок выбрать \texttt{Mann--Whitney U}, для парных разностей выбрать \texttt{Wilcoxon signed-rank}. 3. Через \texttt{ANOVA} выбрать \texttt{One-Way ANOVA}.\par Формулы: $t=(\bar{x}-\mu_0)/(s/\sqrt{n})$ для одного среднего, $F=MS_{\text{между}}/MS_{\text{внутри}}$ для ANOVA. Смысл критериев состоит в сравнении наблюдаемой разности с её стандартной ошибкой.
\end{jamovi}
Выбор критерия определяется числом групп и видом измерений. Непараметрические критерии работают с рангами и подходят при скошенном распределении. Парные критерии требуют, чтобы измерения были связаны одной и той же единицей наблюдения.

\subsection*{Слайд 32. Корреляция и регрессия}
\begin{onslide}
Слайд-разделитель с названием раздела.
\end{onslide}
На слайде показан разделитель раздела «Корреляция и регрессия». Ведущий сообщает, что после сравнения групп рассматриваются связи между непрерывными переменными. Разделитель меняет вопрос с различий на связь. Следующий слайд демонстрирует корреляцию возраста и максимальной ЧСС.

\subsection*{Слайд 33. Корреляция возраста и максимальной ЧСС равна −0,39}
\begin{onslide}
Блок «Выполнено в Jamovi» сообщает, что открыты \texttt{Regression} и \texttt{Correlation Matrix}, переменные \texttt{thalach}, \texttt{age} и \texttt{chol} включены в анализ, рассчитаны корреляции Пирсона и Спирмена. Ниже указано: при $n = 303$ получено $r = -0{,}39$, $p < 0{,}001$, связь умеренная, направление отрицательное. Справа показан фрагмент таблицы корреляции.
\end{onslide}
\begin{jamovi}
Выполнить. 1. Нажать \texttt{Alt}, затем нажать \texttt{A}; \texttt{Tab} до \texttt{Regression}, \texttt{Enter}; выбрать \texttt{Correlation Matrix}, \texttt{Enter}. 2. \texttt{Tab} к списку переменных, выделить \texttt{thalach}, \texttt{age}, \texttt{chol} и перенести нажатием \texttt{Enter}. 3. \texttt{Tab} к опциям, \texttt{Space} включает Pearson и Spearman.
\par Формула: $r=\frac{\sum(x_i-\bar{x})(y_i-\bar{y})}{\sqrt{\sum(x_i-\bar{x})^2\sum(y_i-\bar{y})^2}}$. Смысл: знак коэффициента задаёт направление связи, а модуль описывает её силу.
\end{jamovi}
На слайде в блоке «Выполнено в Jamovi» сообщается, что открыты \texttt{Regression} и \texttt{Correlation Matrix}, переменные \texttt{thalach}, \texttt{age} и \texttt{chol} включены в анализ и рассчитаны корреляции Пирсона и Спирмена. При $n = 303$ получено $r = -0{,}39$, $p < 0{,}001$; связь умеренная, направление отрицательное. Справа показан фрагмент таблицы корреляции. Ведущий указывает на строку таблицы с коэффициентом $-0{,}39$. Корреляция Пирсона оценивает линейную связь, а корреляция Спирмена оценивает монотонную ранговую связь. Следующий слайд объясняет интерпретацию коэффициента.

\subsection*{Слайд 34. Как интерпретировать коэффициент корреляции}
\begin{onslide}
На слайде указано, что коэффициент $r$ принимает значения от $-1$ до $+1$; знак определяет направление связи, модуль характеризует её силу. Для \texttt{thalach} и \texttt{age} коэффициент равен $-0{,}39$; в данном наборе данных связь отрицательная. Корреляция Пирсона оценивает линейную связь и чувствительна к выбросам. Корреляция Спирмена оценивает монотонную ранговую связь и применяется для порядковых данных.
\end{onslide}
На слайде указано, что коэффициент $r$ принимает значения от $-1$ до $+1$, знак определяет направление связи, а модуль характеризует её силу. Для \texttt{thalach} и \texttt{age} коэффициент равен $-0{,}39$, и в данном наборе данных связь отрицательная. Корреляция Пирсона оценивает линейную связь и чувствительна к выбросам. Корреляция Спирмена оценивает монотонную ранговую связь. Ведущий использует формулировку «Выявлена отрицательная корреляция возраста с максимальной ЧСС ($r = -0{,}39$, $p < 0{,}001$)». Корреляция не устанавливает причинность, поскольку общая причина может создавать связь между переменными. Следующий слайд переходит от силы связи к регрессионной модели.

\subsection*{Слайд 35. Линейная регрессия связывает возраст и максимальную ЧСС}
\begin{onslide}
Приведена таблица коэффициентов: Intercept $B = 207{,}97$, $SE = 7{,}78$, $t = 26{,}72$, $p < 0{,}001$, 95\% CI [192,65; 223,29]; \texttt{age} $B = -1{,}02$, $SE = 0{,}13$, $t = -7{,}58$, $p < 0{,}001$, 95\% CI [$-1{,}28$; $-0{,}75$]; \texttt{sex} $B = -4{,}30$, $SE = 2{,}60$, $t = -1{,}66$, $p = 0{,}099$, 95\% CI [$-9{,}42$; 0,81]. Блок сообщает, что $R^2 = 0{,}16$ и что при одинаковом значении пола увеличение возраста на один год ассоциировано со снижением ожидаемой максимальной ЧСС примерно на 1,02 удара в минуту.
\end{onslide}
\begin{jamovi}
Выполнить. 1. Нажать \texttt{Alt}, затем нажать \texttt{A}; \texttt{Tab} до \texttt{Regression}, \texttt{Enter}; выбрать \texttt{Linear Regression}, \texttt{Enter}. 2. \texttt{Tab} к Dependent Variable: \texttt{thalach}, \texttt{Enter}; к Covariate: \texttt{age}, \texttt{Enter}; к Factor: \texttt{sex}, \texttt{Enter}.
\par Формулы: $B=\frac{\sum(x_i-\bar{x})(y_i-\bar{y})}{\sum(x_i-\bar{x})^2}$, $A=\bar{y}-B\bar{x}$, $R^2=1-SSE/SST$. Смысл: коэффициент $B$ показывает изменение отклика при увеличении предиктора на единицу, а $R^2$описывает долю объяснённого разброса.
\end{jamovi}
На слайде приведена таблица коэффициентов: переменная Intercept имеет $B = 207{,}97$, $SE = 7{,}78$, $t = 26{,}72$, $p < 0{,}001$ и 95\% CI [192,65; 223,29]. Переменная \texttt{age} имеет $B = -1{,}02$, $SE = 0{,}13$, $t = -7{,}58$, $p < 0{,}001$ и 95\% CI [$-1{,}28$; $-0{,}75$]. Переменная \texttt{sex} имеет $B = -4{,}30$, $SE = 2{,}60$, $t = -1{,}66$, $p = 0{,}099$ и 95\% CI [$-9{,}42$; 0,81]. Блок сообщает, что $R^2 = 0{,}16$ и что при одинаковом значении пола увеличение возраста на один год ассоциировано со снижением ожидаемой максимальной ЧСС примерно на 1,02 удара в минуту. Ведущий разбирает таблицу: $B$ показывает изменение $Y$ при изменении $X$ на единицу, $p$ оценивает вклад предиктора, $R^2$ описывает долю объяснённого разброса, а доверительный интервал приведён в той же строке. Следующий слайд показывает, как записать результат в статье.

\subsection*{Слайд 36. Как записать результат регрессии в статье}
\begin{onslide}
Приведена формулировка: «При одинаковом значении пола увеличение возраста на один год ассоциировано со снижением ожидаемой максимальной ЧСС приблизительно на 1,02 уд/мин: $B = -1{,}02$ (95\% CI $-1{,}28$; $-0{,}75$), $p < 0{,}001$».
\end{onslide}
На слайде приведена формулировка «При одинаковом значении пола увеличение возраста на один год ассоциировано со снижением ожидаемой максимальной ЧСС приблизительно на 1,02 уд/мин: $B = -1{,}02$ (95\% CI $-1{,}28$; $-0{,}75$), $p < 0{,}001$». Ведущий сообщает, что причинные выводы требуют соответствующего дизайна исследования. Запись переносит числа из таблицы в текст статьи без изменений смысла. Следующий раздел рассматривает категориальные данные.

\subsection*{Слайд 37. Множественная регрессия и конфайндинг}
\begin{onslide}
Модель имеет вид $Y=\beta_0+\beta_1X_1+\beta_2X_2+\varepsilon$. Коэффициент $\beta_1$ описывает связь $X_1$ с $Y$ при фиксированном значении $X_2$. Конфайндинг возникает, когда переменная является общей причиной экспозиции и исхода. Изменение коэффициента при добавлении переменной служит диагностическим признаком, но само по себе не доказывает конфайндинг.
\end{onslide}
\begin{jamovi}
Выполнить. В \texttt{Linear Regression} перенести несколько переменных в поле \texttt{Covariates}.\par Формула модели содержит несколько предикторов и случайную ошибку. Смысл коэффициента состоит в оценке вклада предиктора при неизменных значениях остальных переменных.
\end{jamovi}
Сравнение моделей до и после добавления переменной может указывать на смещение, но не доказывает конфайндинг. Механическое включение всех переменных может привести к контролю медиатора или коллайдера. Вывод зависит от того, какие переменные измерены и включены в анализ.

\subsection*{Слайд 38. Проверки регрессии и интервал корреляции}
\begin{onslide}
На слайде описаны гомоскедастичность и график остатков. Доверительный интервал для коэффициента корреляции строят через преобразование Фишера $z=\frac{1}{2}\ln\frac{1+r}{1-r}$.
\end{onslide}
\begin{jamovi}
Выполнить. 1. В \texttt{Linear Regression} включить график остатков. 2. Для корреляции включить доверительный интервал.\par Формула преобразования Фишера связывает $r$ и $z$. Смысл интервала состоит в диапазоне значений коэффициента, согласующихся с данными.
\end{jamovi}
Гомоскедастичность означает постоянство разброса остатков во всём диапазоне предиктора. Воронка на графике остатков указывает на нарушение этого условия. Интервал для $r$ делает оценку связи более информативной, чем одно точечное значение.

\subsection*{Слайд 39. Категориальные данные}
\begin{onslide}
Слайд-разделитель с названием раздела.
\end{onslide}
На слайде показан разделитель раздела «Категориальные данные». Ведущий сообщает, что категориальные признаки обрабатываются через таблицы сопряжённости. Разделитель подготавливает новый вид таблицы. Следующий слайд показывает таблицу для \texttt{cp} и \texttt{target}.

\subsection*{Слайд 40. Таблица сопряжённости связывает cp и target}
\begin{onslide}
Приведена таблица критерия: Pearson $\chi^2 = 81{,}82$, $df = 3$, $p < 0{,}001$. Блок сообщает, что минимальная ожидаемая частота равна 10,55 и превышает 5, $V$ Крамера равен 0,52, что указывает на умеренную связь типа боли с наличием болезни.
\end{onslide}
\begin{jamovi}
Выполнить. 1. Нажать \texttt{Alt}, затем нажать \texttt{A}; \texttt{Tab} до \texttt{Frequencies}, \texttt{Enter}; выбрать \texttt{Contingency Tables}, затем \texttt{Independent Samples}, \texttt{Enter}. 2. \texttt{Tab} к Rows: \texttt{cp}, \texttt{Enter}; к Columns: \texttt{target}, \texttt{Enter}. 3. \texttt{Tab} к опциям, \texttt{Space} включает $\chi^2$, \texttt{Expected counts} и $V$ Крамера.
\par Формулы: $\chi^2=\sum\frac{(O-E)^2}{E}$, $V=\sqrt{\frac{\chi^2}{n\min(r-1,c-1)}}$. Смысл: хи-квадрат оценивает расхождение наблюдаемых частот с ожидаемыми, а $V$ Крамера измеряет силу связи между категориями.
\end{jamovi}
На слайде приведена таблица критерия: Pearson $\chi^2 = 81{,}82$, $df = 3$, $p < 0{,}001$. Блок сообщает, что минимальная ожидаемая частота равна 10,55 и превышает 5, а $V$ Крамера равен 0,52, что указывает на умеренную связь типа боли с наличием болезни. В Jamovi открываются \texttt{Frequencies}, \texttt{Contingency Tables} и \texttt{Independent Samples}, после чего \texttt{cp} задаётся в строках, а \texttt{target} задаётся в столбцах. В опциях включаются $\chi^2$, \texttt{Expected counts} и $V$ Крамера. Ведущий показывает ожидаемые частоты и значение $V$ Крамера. Следующий слайд объясняет правила интерпретации.

\subsection*{Слайд 41. Как интерпретировать критерий хи-квадрат}
\begin{onslide}
На слайде указано, что статистика $\chi^2$, числа степеней свободы и $p$-значение интерпретируются по тем же правилам, что и результат $t$-критерия; при ожидаемых частотах менее 5 вывод делается осторожно; $V$ Крамера измеряет силу связи между категориальными признаками и принимает значения от 0 до 1.
\end{onslide}
На слайде указано, что статистика $\chi^2$, числа степеней свободы и $p$-значение интерпретируются по тем же правилам, что и результат $t$-критерия. При ожидаемых частотах менее 5 вывод делается осторожно. $V$ Крамера измеряет силу связи между категориальными признаками и принимает значения от 0 до 1. Ведущий сообщает, что изменяется только вид таблицы, а $p$-значение, значимость и величина эффекта сохраняют тот же смысл. Следующий раздел посвящён графикам.

\subsection*{Слайд 42. Точный критерий и логистическая регрессия}
\begin{onslide}
Точный критерий Фишера применяется при малых ожидаемых частотах. Логистическая модель имеет вид $\ln\frac{p}{1-p}=\beta_0+\beta_1X$. Экспонента коэффициента даёт отношение шансов $OR=\exp(\beta_1)$.
\end{onslide}
\begin{jamovi}
Выполнить. 1. В \texttt{Contingency Tables} включить \texttt{Fisher exact}. 2. Через \texttt{Logistic Regression} задать бинарный исход и предиктор.\par Формула логистической модели связывает логарифм шансов с предикторами. Смысл $OR$ состоит в изменении шансов исхода при увеличении предиктора на единицу.
\end{jamovi}
Точный критерий Фишера рассчитывает вероятность наблюдаемой и более крайних таблиц при фиксированных маргинальных суммах. Логистическая регрессия оценивает вероятность события, а не среднее значение непрерывной величины. Отношение шансов требует осторожной интерпретации при высокой распространённости исхода.

\subsection*{Слайд 43. Графики}
\begin{onslide}
Слайд-разделитель с названием раздела.
\end{onslide}
На слайде показан разделитель раздела «Графики». Ведущий сообщает, что графики строятся из тех же переменных и требуют подписей. Разделитель подготавливает визуальное представление результатов. Следующий слайд показывает три базовых графика.

\subsection*{Слайд 44. Ящичная диаграмма: группы и разброс}
\begin{onslide}
На слайде показана ящичная диаграмма распределения \texttt{chol} по \texttt{target}. Ось $X$ показывает группы отсутствия и наличия болезни, а ось $Y$ показывает уровень холестерина в мг/дл.
\end{onslide}
\begin{jamovi}
Выполнить. В \texttt{Descriptives} открыть \texttt{Plots} и включить \texttt{Box plot} для \texttt{chol} по \texttt{target}.\par Обозначения: медиана делит выборку пополам, ящик показывает $IQR=Q_3-Q_1$, а отдельные точки показывают выбросы.
\end{jamovi}
Ведущий объясняет, что ящик показывает основной массив значений, а линия внутри ящика задаёт медиану. Усы показывают границы разброса, а точки за пределами усов показывают выбросы. Сравнение двух ящиков позволяет увидеть различие групп по уровню холестерина. Следующий слайд разбирает диаграмму рассеяния.

\subsection*{Слайд 45. Диаграмма рассеяния: возраст и максимальная ЧСС}
\begin{onslide}
На слайде показана диаграмма рассеяния \texttt{age} и \texttt{thalach}. Ось $X$ показывает возраст в годах, а ось $Y$ показывает максимальную частоту сердечных сокращений в ударах в минуту.
\end{onslide}
\begin{jamovi}
Выполнить. В \texttt{Plots} выбрать \texttt{Scatter Plot} для \texttt{age} и \texttt{thalach}.\par Обозначения: каждая точка соответствует одному пациенту, а линия регрессии показывает направление связи.
\end{jamovi}
Диаграмма рассеяния показывает форму связи и наличие выбросов. Линия регрессии описывает направление изменения максимальной частоты с возрастом. Разброс точек вокруг линии показывает, насколько тесна связь. Следующий слайд переходит к гистограмме.

\subsection*{Слайд 46. Гистограмма: распределение возраста}
\begin{onslide}
На слайде показана гистограмма возраста. Ось $X$ показывает возраст в годах, а ось $Y$ показывает число пациентов.
\end{onslide}
\begin{jamovi}
Выполнить. Гистограмма построена в блоке нормальности через \texttt{Descriptives} и \texttt{Plots}.\par Обозначения: высота столбца задаёт частоту наблюдений в интервале.
\end{jamovi}
Гистограмма показывает форму распределения и позволяет заметить скошенность. В блоке приведены элементы реальных статей, которые покрывают около 95\% типичных случаев, за некоторым исключением редких методов. Ведущий связывает три графика с типами данных и переходит к требованиям к подписям.

\subsection*{Слайд 47. Кривые Каплана–Мейера: время до события}
\begin{onslide}
На слайде показаны схематические кривые Каплана--Мейера. Ось $X$ показывает время после начала наблюдения, а ось $Y$ показывает долю пациентов без события.
\end{onslide}
Каждый уступ кривой соответствует событию, а вертикальные отметки показывают цензурирование. Медиана времени до события считывается по пересечению кривой с уровнем 0,5. Ведущий объясняет, что число пациентов под риском приведено под осью и что различие кривых дополняется отношением рисков с доверительным интервалом.

\subsection*{Слайд 48. Форест-плот: эффекты исследований}
\begin{onslide}
На слайде показана схема форест-плота. Ось $X$ показывает величину эффекта, а каждая строка соответствует исследованию.
\end{onslide}
Точка задаёт оценку эффекта, а линия показывает 95\% доверительный интервал. Вертикальная линия нуля отделяет отсутствие эффекта, а ромб показывает объединённую оценку. Ведущий указывает, что пересечение интервала с нулём указывает на отсутствие доказательства эффекта.

\subsection*{Слайд 49. График остатков: проверка модели}
\begin{onslide}
На слайде показан график остатков. Ось $X$ показывает предсказанные значения, а ось $Y$ показывает остатки.
\end{onslide}
Точки должны располагаться вокруг нуля без воронки и систематических дуг. Воронка указывает на нарушение гомоскедастичности. Ведущий связывает график с проверкой линейной регрессии, рассмотренной ранее. В демонстрациях используется около десяти тестов, поэтому выводы следует рассматривать с оговоркой о множественности.

\subsection*{Слайд 50. ROC-кривая: точность бинарного классификатора}
\begin{onslide}
На слайде показана ROC-кривая. Ось $X$ показывает долю ложноположительных результатов, а ось $Y$ показывает долю истинноположительных результатов.
\end{onslide}
Каждая точка кривой соответствует порогу принятия решения. Площадь под кривой описывает качество классификации, а пунктирная диагональ соответствует случайному угадыванию. Ведущий объясняет, что кривая применяется к бинарным исходам, например к наличию болезни.

\subsection*{Слайд 51. Подпись графика содержит обязательные элементы}
\begin{onslide}
На слайде указано, что подпись содержит оси $X$ и $Y$, заголовок и единицы измерения; в Jamovi эти элементы задаются в разделе \texttt{Plot \& Axis Titles}.
\end{onslide}
На слайде указано, что подпись содержит оси $X$ и $Y$, заголовок и единицы измерения. В Jamovi эти элементы задаются в разделе \texttt{Plot \& Axis Titles}. Ведущий отмечает типичную ошибку, когда на осях отсутствуют единицы измерения. Подпись делает график понятным вне контекста презентации. Следующий раздел знакомит с байесовским статистическим выводом.

\subsection*{Слайд 52. Байесовский статистический вывод}
\begin{onslide}
Слайд-разделитель с названием раздела.
\end{onslide}
На слайде показан разделитель раздела «Байесовский статистический вывод». Ведущий сообщает, что вторым способом вывода о гипотезах является байесовский статистический вывод. Разделитель вводит альтернативу частотному подходу. Следующий слайд сравнивает вопросы, которые задают два подхода.

\subsection*{Слайд 53. Частотный критерий и байесовский статистический вывод задают разные во-}
\begin{onslide}
В левом блоке приведён вопрос частотного теста: насколько неожиданны данные, если $H_0$ верна, и формула $p = P(|T| \ge |t_{\text{набл}}| \mid H_0)$. В правом блоке приведён вопрос байесовского вывода: во сколько раз данные согласуются с $H_1$ лучше, чем с $H_0$, и формула $\mathrm{BF}_{10} = p(D\mid H_1)/p(D\mid H_0)$. Ниже указано, что $p$-значение описывает данные при заданной гипотезе, а байесовский фактор обновляет априорные шансы гипотез. Вертикальная черта в формуле читается как «при условии». Символ $H_0$ обозначает нулевую гипотезу об отсутствии различия. Слова «данные или более крайнее» означают наблюдаемый результат и результаты с таким же или большим различием. Тогда $p$-значение равно вероятности получить такой или более крайний результат при условии, что нулевая гипотеза верна.
\end{onslide}
На слайде в левом блоке приведён вопрос частотного теста: насколько неожиданны данные, если $H_0$ верна, с формулой $p = P(|T| \ge |t_{\text{набл}}| \mid H_0)$. В правом блоке приведён вопрос байесовского вывода: во сколько раз данные согласуются с $H_1$ лучше, чем с $H_0$, с формулой $\mathrm{BF}_{10} = p(D\mid H_1)/p(D\mid H_0)$. Ниже указано, что $p$-значение описывает данные при заданной гипотезе, а байесовский фактор обновляет априорные шансы гипотез. Ведущий поясняет, что вертикальная черта в формуле читается как «при условии», а символ $H_0$ обозначает нулевую гипотезу об отсутствии различия. Условие $|T| \ge |t_{\text{набл}}|$ означает наблюдаемое значение статистики и все более крайние значения. Тогда $p$-значение равно вероятности получить такой или более крайний результат при условии, что нулевая гипотеза верна; обоснование приведено у Greenland et al. (2016) и Wasserstein \& Lazar (2016), а следующий слайд показывает переход от априора к апостериору.

\subsection*{Слайд 54. От априорного распределения к апостериорному}
\begin{onslide}
На рисунке показаны три кривые. Три шага: априорное распределение описывает представления до наблюдения данных; функция правдоподобия отражает информацию в данных; апостериорное распределение описывает обновлённые представления. Под рисунком в левой колонке приведена формула $p(\theta \mid D) \propto p(D \mid \theta)\,p(\theta)$ с обозначениями: $p(\theta \mid D)$ — апостериорное распределение; $p(D \mid \theta)$ — функция правдоподобия; $p(\theta)$ — априорное распределение. В правой колонке расположен блок «Три этапа вывода».
\end{onslide}
На рисунке показаны три кривые. Априорное распределение описывает представления до наблюдения данных, функция правдоподобия отражает информацию в данных, а апостериорное распределение описывает обновлённые представления. Под рисунком приведена формула $p(\theta \mid D) \propto p(D \mid \theta)\,p(\theta)$ с обозначениями: $p(\theta \mid D)$ — апостериорное распределение, $p(D \mid \theta)$ — функция правдоподобия, $p(\theta)$ — априорное распределение. В правой колонке расположен блок «Три этапа вывода». Ведущий сообщает, что при большом объёме данных влияние априорного распределения уменьшается. Литература указана у Wagenmakers et al. (2018, Part I) и Kruschke (2013), а следующий слайд вводит шкалу байесовского фактора.

\subsection*{Слайд 55. Шкала байесовского фактора}
\begin{onslide}
На рисунке приведена шкала: $\mathrm{BF}_{10} < 1/3$ указывает на данные в пользу $H_0$, значения от $1/3$ до 1 соответствуют слабым доказательствам в пользу $H_0$, значения от 1 до 3 соответствуют слабым доказательствам в пользу $H_1$, значения от 3 до 10 соответствуют умеренным, от 10 до 30 соответствуют сильным, значения более 30 соответствуют очень сильным доказательствам. Ниже указано, что при $\mathrm{BF}_{10} = 10$ данные в 10 раз более согласованы с $H_1$, чем с $H_0$, а при $\mathrm{BF} \approx 1$ данных недостаточно.
\end{onslide}
На рисунке приведена шкала байесовского фактора. Значение $\mathrm{BF}_{10} < 1/3$ указывает на данные в пользу $H_0$, значения от $1/3$ до 1 соответствуют слабым доказательствам в пользу $H_0$, а значения от 1 до 3 соответствуют слабым доказательствам в пользу $H_1$. Значения от 3 до 10 соответствуют умеренным доказательствам, значения от 10 до 30 соответствуют сильным доказательствам, а значения более 30 соответствуют очень сильным доказательствам. Ниже указано, что при $\mathrm{BF}_{10} = 10$ данные в 10 раз более согласованы с $H_1$, чем с $H_0$, а при $\mathrm{BF} \approx 1$ данных недостаточно. Источником являются Jeffreys (1961) и Lee \& Wagenmakers (2013), а байесовский фактор представляет собой отношение правдоподобий. Следующий слайд применяет шкалу к данным датасета.

\subsection*{Слайд 56. Байесовский фактор на данных датасета}
\begin{onslide}
Таблица содержит строки: \texttt{chol} по \texttt{sex}: $t = 3{,}54$, $p < 0{,}001$, $\mathrm{BF}_{10} = 26{,}5$; \texttt{thalach} по \texttt{sex}: $t = -0{,}85$, $p = 0{,}399$, $\mathrm{BF}_{10} = 0{,}22$. Ниже указано, что для \texttt{chol} получены сильные доказательства в пользу различий, для \texttt{thalach} $\mathrm{BF}_{01} = 4{,}5$ указывает на умеренные доказательства отсутствия различий. Направление разности — \texttt{sex=0} минус \texttt{sex=1}. Значения BF воспроизведены скриптом \texttt{scripts/verify\_bf.py} для объединённой дисперсии и априора JZS с $r = 0{,}707$.
\end{onslide}
\begin{jamovi}
Выполнить. 1. Открыть \texttt{Independent Samples T-Test}, как на слайде 25. 2. \texttt{Tab} к опции \texttt{Bayes factor}, \texttt{Space} включает её. 3. Повторить для \texttt{chol} и \texttt{thalach}.
\par Формула: $\mathrm{BF}_{10}=\frac{P(D\mid H_1)}{P(D\mid H_0)}$. Смысл: байесовский фактор показывает, во сколько раз данные более согласованы с $H_1$, чем с $H_0$.
\end{jamovi}
В таблице приведены результаты для двух переменных: для \texttt{chol} по \texttt{sex} получены $t = 3{,}54$, $p < 0{,}001$ и $\mathrm{BF}_{10} = 26{,}5$. Для \texttt{thalach} по \texttt{sex} получены $t = -0{,}85$, $p = 0{,}399$ и $\mathrm{BF}_{10} = 0{,}22$. Ниже указано, что для \texttt{chol} получены сильные доказательства в пользу различий, а для \texttt{thalach} значение $\mathrm{BF}_{01} = 4{,}5$ указывает на умеренные доказательства отсутствия различий. Направление разности — \texttt{sex=0} минус \texttt{sex=1}. Значения BF воспроизведены скриптом \texttt{scripts/verify\_bf.py} для объединённой дисперсии и априора JZS с $r = 0{,}707$. Ведущий сообщает, что значение $p > 0{,}05$ само по себе не обосновывает вывод об отсутствии эффекта; такое основание даёт $\mathrm{BF}_{01} > 3$. Следующий слайд переходит к оцениванию параметров.

\subsection*{Слайд 57. Байесовское оценивание параметров}
\begin{onslide}
На рисунке показаны апостериорное распределение, апостериорный интервал и область практической эквивалентности. Блок перечисляет: апостериорное распределение показывает правдоподобные значения эффекта; 95\% апостериорный интервал (credible interval, CrI) содержит 95\% апостериорной вероятности; область практической эквивалентности (ROPE) задаёт диапазон пренебрежимо малых значений эффекта.
\end{onslide}
На рисунке показаны апостериорное распределение, апостериорный интервал и область практической эквивалентности. Апостериорное распределение показывает правдоподобные значения эффекта. 95\% апостериорный интервал (credible interval, CrI) содержит 95\% апостериорной вероятности. Область практической эквивалентности (ROPE) задаёт диапазон пренебрежимо малых значений эффекта. Ведущий сообщает, что апостериорный интервал по смыслу отличается от доверительного интервала. Следующий слайд объясняет, как сообщить результат байесовского анализа.

\subsection*{Слайд 58. Как сообщить результат байесовского анализа}
\begin{onslide}
Порядок сообщения: тип теста и априор, JZS $t$-тест с априором Коши $r = 0{,}707$; байесовский фактор $\mathrm{BF}_{10}$ или $\mathrm{BF}_{01}$ с расшифровкой стороны доказательств; для оценивания сообщаются медиана и 95\% апостериорный интервал, а также отношение к ROPE; проверяется чувствительность к априору. Блок «Частая ошибка» сообщает, что байесовский фактор является отношением правдоподобий, а утверждение «$\mathrm{BF}_{10} = 26{,}5$ означает 96\% вероятность $H_1$» ошибочно.
\end{onslide}
На слайде приведён порядок сообщения результата. Сначала указываются тип теста и априор, например JZS $t$-тест с априором Коши $r = 0{,}707$. Затем сообщается байесовский фактор $\mathrm{BF}_{10}$ или $\mathrm{BF}_{01}$ с расшифровкой стороны доказательств. Для оценивания сообщаются медиана и 95\% апостериорный интервал, а также отношение к ROPE, и проверяется чувствительность к априору. Блок «Частая ошибка» сообщает, что байесовский фактор является отношением правдоподобий, а утверждение «$\mathrm{BF}_{10} = 26{,}5$ означает 96\% вероятность $H_1$» ошибочно. Ведущий добавляет, что отчёт без указания априора и анализа чувствительности невоспроизводим, а следующий слайд показывает формулу Байеса в форме шансов.

\subsection*{Слайд 59. Теорема Байеса в форме шансов}
\begin{onslide}
На рисунке показана схема. Формула: апостериорные шансы равны априорным шансам, умноженным на $\mathrm{BF}_{10}$. Ниже указано, что при априорных шансах 1:1 апостериорные шансы совпадают с байесовским фактором; при других априорных шансах байесовский фактор не равен итоговой вероятности гипотезы.
\end{onslide}
На рисунке показана схема. Формула гласит, что апостериорные шансы равны априорным шансам, умноженным на $\mathrm{BF}_{10}$. При априорных шансах 1:1 апостериорные шансы совпадают с байесовским фактором. При других априорных шансах байесовский фактор не равен итоговой вероятности гипотезы. Ведущий указывает на формулу и поясняет роль каждого множителя. Следующий слайд иллюстрирует роль априорной информации на примере диагностического теста.

\subsection*{Слайд 60. Диагностический тест как априорная информация}
\begin{onslide}
На схеме приведены натуральные частоты: 10\,000 человек, распространённость 1\%, чувствительность 90\%, специфичность 95\%; 90 истинноположительных и 495 ложноположительных результатов. Прогностическая ценность положительного результата (PPV) составляет 15,4\%. Ниже указано, что априорное распределение задаёт исходные шансы гипотезы, а данные их обновляют.
\end{onslide}
На схеме приведены натуральные частоты: 10\,000 человек, распространённость 1\%, чувствительность 90\% и специфичность 95\%. Получены 90 истинноположительных и 495 ложноположительных результатов. Прогностическая ценность положительного результата (PPV) составляет 15,4\%. Ниже указано, что априорное распределение задаёт исходные шансы гипотезы, а данные их обновляют. Ведущий указывает на расчёт $90/(90+495) = 15{,}4\%$. Следующий слайд проверяет, насколько вывод зависит от выбора априора.

\subsection*{Слайд 61. Чувствительность байесовского фактора к априору}
\begin{onslide}
На графике показан байесовский фактор как функция параметра $r$ при $t = 3{,}54$ и $n = 303$. Блок сообщает, что вывод «сильные доказательства» сохраняется при $r$ от 0,3 до 1,5. Ниже указано, что априорное распределение указывается в отчёте и устойчивость вывода проверяется.
\end{onslide}
На графике показан байесовский фактор как функция параметра $r$ при $t = 3{,}54$ и $n = 303$. Блок сообщает, что вывод «сильные доказательства» сохраняется при $r$ от 0,3 до 1,5. Ниже указано, что априорное распределение указывается в отчёте и что устойчивость вывода проверяется. Ведущий указывает диапазон значений $r$, при которых вывод сохраняется. Следующий слайд сравнивает $p$-значение и байесовский фактор при разном объёме выборки.

\subsection*{Слайд 62. p-значение и байесовский фактор при разном объёме выборки}
\begin{onslide}
Таблица содержит строки: $n = 30$, $t = 1{,}10$, $p = 0{,}283$, $\mathrm{BF}_{10} = 0{,}41$; $n = 100$, $t = 2{,}00$, $p = 0{,}048$, $\mathrm{BF}_{10} = 0{,}82$; $n = 300$, $t = 3{,}46$, $p = 0{,}001$, $\mathrm{BF}_{10} = 20{,}9$. Блок сообщает, что при $n = 100$ получено $p < 0{,}05$ и $\mathrm{BF}_{10} = 0{,}82$; данные скорее согласуются с $H_0$.
\end{onslide}
В таблице приведены три варианта расчёта: при $n = 30$ получены $t = 1{,}10$, $p = 0{,}283$ и $\mathrm{BF}_{10} = 0{,}41$. При $n = 100$ получены $t = 2{,}00$, $p = 0{,}048$ и $\mathrm{BF}_{10} = 0{,}82$. При $n = 300$ получены $t = 3{,}46$, $p = 0{,}001$ и $\mathrm{BF}_{10} = 20{,}9$. Блок сообщает, что при $n = 100$ получено $p < 0{,}05$ и $\mathrm{BF}_{10} = 0{,}82$, то есть данные скорее согласуются с $H_0$. Ведущий сообщает, что разные вопросы дают разные ответы. Следующий слайд показывает байесовский фактор при последовательном накоплении данных.

\subsection*{Слайд 63. Байесовский фактор при последовательном накоплении данных}
\begin{onslide}
На графике показан байесовский фактор при добавлении наблюдений для фиксированного эффекта $d = 0{,}4$. Блок сообщает, что при добавлении наблюдений байесовский фактор остаётся отношением правдоподобий и заранее фиксированный объём выборки не требуется. Ниже указано, что планирование объёма выборки всё равно необходимо.
\end{onslide}
На графике показан байесовский фактор при добавлении наблюдений для фиксированного эффекта $d = 0{,}4$. Блок сообщает, что при добавлении наблюдений байесовский фактор остаётся отношением правдоподобий и что заранее фиксированный объём выборки не требуется. Ниже указано, что планирование объёма выборки всё равно необходимо. Ведущий указывает на изменение кривой по мере накопления данных. Следующий слайд знакомит с марковскими цепями Монте-Карло.

\subsection*{Слайд 64. Марковские цепи Монте-Карло (MCMC) и апостериорное распределение}
\begin{onslide}
На рисунке показаны график траектории цепи и гистограмма апостериорной выборки. Блок «Проверка сходимости» содержит три условия: график траектории цепи не имеет тренда и участков застревания; несколько цепочек дают совпадающие распределения; эффективный размер выборки ESS достаточно велик.
\end{onslide}
На рисунке показаны график траектории цепи и гистограмма апостериорной выборки. Блок «Проверка сходимости» содержит три условия: график траектории цепи не имеет тренда и участков застревания, несколько цепочек дают совпадающие распределения, а эффективный размер выборки ESS достаточно велик. Ведущий сообщает, что проверка сходимости обязательна для вычислений на основе марковских цепей Монте-Карло. Следующий слайд перечисляет ограничения байесовского статистического вывода.

\subsection*{Слайд 65. Возможности байесовского вывода}
\begin{onslide}
На слайде перечислены возможности байесовского подхода со ссылками на источники: полное распределение правдоподобных значений параметров, учёт предыдущих находок через априор, оценка доказательств для $H_0$ и $H_1$, отслеживание доказательств по мере накопления данных и правило HDI+ROPE для практического принятия нулевого значения.
\end{onslide}
Пояснение опирается на работы Kruschke (2013, 2018) и Wagenmakers et al. (2018, Part I). Ведущий подчёркивает, что возможности реализуются при заданных априорных распределениях и адекватной модели. Следующий слайд описывает ограничения метода.
\subsection*{Слайд 66. Ограничения байесовского статистического вывода}
\begin{onslide}
На слайде перечислены четыре положения: выбор априорного распределения субъективен и влияет на результат при малом объёме выборки; вычисления на основе марковских цепей Монте-Карло требуют проверки сходимости; байесовский фактор дополняет оценку дизайна, качества данных и клинического смысла; отчёт без указания априора и анализа чувствительности невоспроизводим. Корректная формулировка: «при заданном априоре $\mathrm{BF}_{10} = 26{,}5$; вывод устойчив к изменению априора».
\end{onslide}
На слайде перечислены четыре положения: выбор априорного распределения субъективен и влияет на результат при малом объёме выборки. Вычисления на основе марковских цепей Монте-Карло требуют проверки сходимости. Байесовский фактор дополняет оценку дизайна, качества данных и клинического смысла. Отчёт без указания априора и анализа чувствительности невоспроизводим. Ведущий использует корректную формулировку «при заданном априоре $\mathrm{BF}_{10} = 26{,}5$; вывод устойчив к изменению априора». Локальные файлы литературы перечислены под блоком, а следующий раздел применяет навыки к клинической статье.

\subsection*{Слайд 67. Как читать клиническую статью}
\begin{onslide}
Слайд-разделитель с названием раздела.
\end{onslide}
На слайде показан разделитель раздела «Как читать клиническую статью». Ведущий сообщает, что полученные навыки применяются к реальной публикации. Разделитель переводит изложение от расчётов к критическому чтению. Следующий слайд начинает разбор со статьи Johansson et al., 2022.

\subsection*{Слайд 68. Заголовок и абстракт статьи}
\begin{onslide}
На слайде показан фрагмент первой страницы статьи Johansson et al., 2022. Блок сообщает, что статья опубликована в Annals of Neurology, DOI 10.1002/ana.26291, и что исследование Park-in-Shape является двойным слепым РКИ с участием 130 пациентов. Ниже указано, что структура PICO включает популяцию, вмешательство, сравнение и исход, а абстракт содержит рубрики Objective, Methods, Results и Interpretation.
\end{onslide}
На слайде показан фрагмент первой страницы статьи Johansson et al., 2022. Блок сообщает, что статья опубликована в Annals of Neurology, DOI 10.1002/ana.26291, а исследование Park-in-Shape является двойным слепым РКИ с участием 130 пациентов. Ниже указано, что структура PICO включает популяцию, вмешательство, сравнение и исход, а абстракт содержит рубрики Objective, Methods, Results и Interpretation. Источником служит Johansson ME, Cameron IGM, Van der Kolk NM, et al. Aerobic Exercise Alters Brain Function and Structure in Parkinson's Disease: A Randomized Controlled Trial. Ann Neurol, 2022;91(2):203--216. Локальная копия статьи хранится в файле \texttt{literature/rct-parkinson-exercise.pdf}, а ведущий указывает на выделенные фразы в заголовке и абстракте; следующий слайд разбирает дизайн исследования.

\subsection*{Слайд 69. Дизайн исследования и контрольная группа}
\begin{onslide}
На слайде показан фрагмент раздела методов с выделениями. Перечислено: распределение по группам выполнено 1:1; исследование двойное слепое; контрольную группу составила растяжка; вмешательство длилось 6 месяцев и включало три занятия в неделю; подгруппа МРТ включала 25 и 31 пациента без отбора.
\end{onslide}
На слайде показан фрагмент раздела методов с выделениями. Перечислено, что распределение по группам выполнено 1:1, что исследование двойное слепое и что контрольную группу составила растяжка. Вмешательство длилось 6 месяцев и включало три занятия в неделю. Подгруппа МРТ включала 25 и 31 пациента без отбора. Ведущий указывает на выделенные фразы «1:1 ratio», «double-blind», «active control (stretching)» и «6 months». Следующий слайд показывает поток участников.

\subsection*{Слайд 70. Поток участников по CONSORT}
\begin{onslide}
На рисунке показан поток участников. 130 пациентов участвовали в исходном исследовании Park-in-Shape. Для дополнительного МРТ-исследования отобрали 57 последовательных совместимых со сканером добровольцев; в группах аэробной нагрузки и растяжки оказались 25 и 31 пациент. Ниже указано, что причины потерь связаны со сканированием, отказом пациентов и техническими причинами и что поток участников является обязательным элементом отчёта.
\end{onslide}
На рисунке показан поток участников. 130 пациентов участвовали в исходном исследовании Park-in-Shape. Для дополнительного МРТ-исследования отобрали 57 последовательных совместимых со сканером добровольцев; в группах аэробной нагрузки и растяжки оказались 25 и 31 пациент. Ниже указано, что причины потерь связаны со сканированием, отказом пациентов и техническими причинами. Поток участников является обязательным элементом отчёта. Ведущий указывает на выделенные числа 130, 57, 25 и 31, а следующий слайд переходит к таблице характеристик.

\subsection*{Слайд 71. Таблица 1 описывает характеристики пациентов}
\begin{onslide}
На слайде показан фрагмент таблицы 1. Перечислено: таблица описывает сопоставимость групп; проверяются $n$, возраст, пол и тяжесть; приведены $p$-значения; непрерывные признаки описаны средним и стандартным отклонением; категориальные признаки описаны частотами. Ниже указано, что случайные различия возможны при корректной рандомизации и что тесты значимости исходных различий не используют как проверку качества рандомизации.
\end{onslide}
На слайде показан фрагмент таблицы 1. Таблица описывает сопоставимость групп, и в ней проверяются $n$, возраст, пол и тяжесть. Приведены $p$-значения. Непрерывные признаки описаны средним и стандартным отклонением, а категориальные признаки описаны частотами. Ниже указано, что случайные различия возможны при корректной рандомизации и что тесты значимости исходных различий не используют как проверку качества рандомизации. Ведущий указывает на выделенные строки Participants, Age и Baseline MDS-UPDRS-III, а следующий слайд объясняет чтение ящичной диаграммы.

\subsection*{Слайд 72. Как читать ящичную диаграмму в статье}
\begin{onslide}
На рисунке показан фрагмент статьи с ящичной диаграммой. Перечислено: медиана и ящик показывают межквартильный размах; усы показывают границы разброса, отдельные точки показывают выбросы; подписи групп и единицы измерения оси указываются обязательно; знак $*$ соответствует сноске $p < 0{,}05$.
\end{onslide}
На рисунке показан фрагмент статьи с ящичной диаграммой. Медиана и ящик показывают межквартильный размах. Усы показывают границы разброса по правилам, заданным при построении; отдельные точки показывают значения за пределами усов. Подписи групп и единицы измерения оси указываются обязательно. Знак $*$ соответствует сноске $p < 0{,}05$. Ведущий указывает на выделенные фразы подписи рисунка «Lower and upper whiskers», «outlying data points» и «*p < 0.05», а следующий слайд разбирает корреляцию в статье.

\subsection*{Слайд 73. Корреляция связности и подготовленности: ρ = 0,62, p = 0,003}
\begin{onslide}
На рисунке показан фрагмент рисунка 5C. Ось $X$ показывает изменение связности DLPFC--RFPN, ось $Y$ показывает изменение VO$_{2}$max. Каждая точка соответствует одному пациенту ($n = 21$). Линия демонстрирует аппроксимацию, коэффициент $\rho$ показывает силу связи.
\end{onslide}
На рисунке показан фрагмент рисунка 5C. Ось $X$ показывает изменение связности DLPFC--RFPN, а ось $Y$ показывает изменение VO$_{2}$max. Каждая точка соответствует одному пациенту ($n = 21$). Линия демонстрирует аппроксимацию, а коэффициент $\rho$ показывает силу связи. Ведущий указывает на выделенную подпись «rho = 0.62, p = 0.003», а следующий слайд показывает, как результат расположен в тексте статьи.

\subsection*{Слайд 74. Как результат расположен в тексте статьи}
\begin{onslide}
На рисунке показан фрагмент раздела результатов. Сообщение результата содержит эффект, коэффициент, $p$-значение, объём выборки и ссылку на рисунок. В статье указаны $\rho = 0{,}62$, $p = 0{,}003$, $n = 21$ и ссылка на рисунок 5C.
\end{onslide}
На рисунке показан фрагмент раздела результатов. Сообщение результата содержит эффект, коэффициент, $p$-значение, объём выборки и ссылку на рисунок. В статье указаны $\rho = 0{,}62$, $p = 0{,}003$, $n = 21$ и ссылка на рисунок 5C. Ведущий указывает на выделенный фрагмент с числами, а следующий слайд рассматривает ограничения в разделе обсуждения.

\subsection*{Слайд 75. Ограничения в разделе обсуждения}
\begin{onslide}
На рисунках показаны выделенные фразы из раздела обсуждения. Один блок пояснений сообщает, что подгруппа МРТ формировалась без случайного отбора, анализы носят поисковый характер, имеются множественные сравнения, а выводы ограничены ранней стадией болезни Паркинсона.
\end{onslide}
На рисунках показаны выделенные фразы из раздела обсуждения. Один блок пояснений сообщает, что подгруппа МРТ формировалась без случайного отбора. Анализы носят поисковый характер, имеются множественные сравнения, а выводы ограничены ранней стадией болезни Паркинсона. Ведущий указывает на выделенную надпись «Limitations of This Study» и на фразу «multiple comparisons», а следующий слайд объясняет, как дизайн определяет силу вывода.

\subsection*{Слайд 76. Дизайн исследования определяет силу вывода}
\begin{onslide}
На рисунке приведена пирамида дизайнов. Перечислено: исследование случая формулирует гипотезу; исследование случай--контроль оценивает ассоциацию; когортное исследование описывает динамику и частоту событий; РКИ обеспечивает контроль факторов; систематический обзор синтезирует результаты исследований. Ниже указано, что дизайн исследования определяет допустимый вывод и статистический анализ не исправляет дизайн.
\end{onslide}
На рисунке приведена пирамида дизайнов. Исследование случая формулирует гипотезу, исследование случай--контроль оценивает ассоциацию, а когортное исследование описывает динамику и частоту событий. РКИ обеспечивает контроль факторов, а систематический обзор синтезирует результаты исследований. Ниже указано, что дизайн исследования определяет допустимый вывод и что статистический анализ не исправляет дизайн. Ведущий указывает на пирамиду снизу вверх, а следующий слайд раскрывает требования к рандомизации.
Иерархия зависит от вопроса. Тип дизайна не гарантирует достоверность результата.

\subsection*{Слайд 77. Рандомизация и сокрытие распределения}
\begin{onslide}
На рисунке показан фрагмент раздела методов. В этой публикации указаны распределение 1:1, ослепление участников в отношении альтернативного вмешательства и ослепление оценщиков. Сведений о генерации последовательности, сокрытии распределения, блоковой и стратифицированной рандомизации здесь нет.
\end{onslide}
На рисунке показан фрагмент раздела методов. Ведущий указывает на выделенные фразы «1:1 ratio», «double-blind» и «active control». В публикации нет сведений о генерации последовательности и сокрытии распределения; такие сведения могут находиться в статье основного исследования. Ослепление врача не упомянуто. Следующий слайд разбирает ITT-анализ и потери участников.

\subsection*{Слайд 78. ITT-анализ и потери участников}
\begin{onslide}
На рисунке показан поток участников. Сообщается, что ITT-анализ включает всех рандомизированных пациентов в их исходных группах. Анализ по протоколу исключает нарушителей и может сместить оценку эффекта. В статье приведены причины выбытия по каждому этапу; отсутствие таких сведений является поводом для критики.
\end{onslide}
На рисунке показан поток участников. Сообщается, что ITT-анализ включает всех рандомизированных пациентов в их исходных группах. Анализ по протоколу исключает нарушителей и может сместить оценку эффекта. В статье приведены причины выбытия по каждому этапу, а отсутствие таких сведений является поводом для критики. Ведущий указывает на блоки с причинами выбытия и на число завершивших наблюдение, а следующий слайд описывает статистический план статьи.

\subsection*{Слайд 79. Статистический план статьи}
\begin{onslide}
На рисунке показан фрагмент раздела статистики. Перечислено: описаны первичная конечная точка и модель анализа, например ANCOVA или смешанные эффекты; проведена корректировка на множественные сравнения; указана совокупность пациентов, включённых в анализ; описан способ обработки пропущенных значений.
\end{onslide}
На рисунке показан фрагмент раздела статистики. Описаны первичная конечная точка и модель анализа, например ANCOVA или смешанные эффекты. Проведена корректировка на множественные сравнения. Указана совокупность пациентов, включённых в анализ. Описан способ обработки пропущенных значений. Ведущий указывает на выделенные фразы «Multiple Comparisons», «familywise error corrections» и «exploratory», а следующий слайд показывает таблицу с величиной эффекта.

\subsection*{Слайд 80. Таблица 3 содержит эффект, g Хеджеса и доверительный интервал}
\begin{onslide}
На рисунке показана таблица 3. Перечислено: анатомическая область и размер кластера описаны в таблице; $p$-значение скорректировано на ошибку первого рода для семейства гипотез (FWER); величина эффекта $g$ Хеджеса равна 1,29 при 95\% доверительном интервале [0,65; 1,91].
\end{onslide}
На рисунке показана таблица 3. В ней описаны анатомическая область и размер кластера. $p$-значение скорректировано на ошибку первого рода для семейства гипотез (FWER). Величина эффекта $g$ Хеджеса равна 1,29 при 95\% доверительном интервале [0,65; 1,91]. Ведущий указывает на выделенный заголовок «Hedges' g (95\% CI)», значение $p = 0{,}027$ и значение 1,29 [0,65; 1,91], а следующий слайд разбирает подпись к ящичной диаграмме.

\subsection*{Слайд 81. Рисунок 4: подпись к ящичной диаграмме}
\begin{onslide}
На рисунке показана подпись к рисунку 4. Сноски $^*p < 0{,}05$ и $^{***}p < 0{,}001$ задают пороги значимости.
\end{onslide}
На рисунке показана подпись к рисунку 4. Сноски $*p < 0{,}05$ и $***p < 0{,}001$ задают пороги значимости. Ведущий указывает на выделенную сноску с тремя звёздочками. Пороги позволяют правильно считать обозначения на графике. Следующий слайд вводит форест-плот.

\subsection*{Слайд 82. Форест-плот показывает эффекты исследований}
\begin{onslide}
На рисунке показана схема форест-плота. Точка показывает оценку эффекта, а линия показывает 95\% доверительный интервал. Вертикальная линия задаёт отсутствие эффекта: 0 для разностных мер и 1 для относительных мер: OR, HR и относительного риска RR. Интервал, включающий это значение, не исключает отсутствие эффекта.
\end{onslide}
На рисунке показана схема форест-плота. Точка показывает оценку эффекта, а линия показывает 95\% доверительный интервал. Вертикальная линия задаёт отсутствие эффекта: 0 для разностных мер и 1 для относительных мер: OR, HR и относительного риска RR. Интервал, включающий это значение, не исключает отсутствие эффекта. Ромб показывает объединённую оценку. Ведущий указывает на точки, линии, вертикальную линию нуля и ромб, а следующий слайд переходит к кривым Каплана--Мейера.

\subsection*{Слайд 83. Оценка Каплана–Мейера описывает время до события}
\begin{onslide}
На рисунке показаны схематические кривые. Медиана времени до события показывает центральное значение. Число пациентов под риском приведено под осью. Вертикальные отметки обозначают цензурирование. Различие кривых дополняется отношением рисков (HR) с 95\% доверительным интервалом.
\end{onslide}
На рисунке показаны схематические кривые. Медиана времени до события показывает центральное значение. Число пациентов под риском приведено под осью. Вертикальные отметки обозначают цензурирование. Различие кривых дополняется отношением рисков (HR) с 95\% доверительным интервалом. Ведущий указывает на медиану, строку «под риском» и вертикальные отметки, а следующий слайд переводит статистический результат в клиническое решение.

\subsection*{Слайд 84. От статистического результата к клиническому решению}
\begin{onslide}
Перечислено: OR описывает отношение шансов, HR описывает отношение мгновенных интенсивностей события среди лиц под риском. ARR равен разнице абсолютных рисков с указанием направления и горизонта. NNT равен 1/ARR и относится к конкретному сравнению, исходу и периоду. NNH рассчитывают для вредного исхода. Ниже указано, что клиническая важность определяется абсолютным эффектом и NNT.
\end{onslide}
На слайде перечислены клинические показатели. OR описывает отношение шансов, а HR описывает отношение мгновенных интенсивностей события в модели пропорциональных рисков. ARR равен разнице абсолютных рисков между контролем и лечением на заданном горизонте. NNT равен 1/ARR и относится к этому сравнению. NNH рассчитывают для одного вредного исхода. Ниже указано, что клиническая важность определяется абсолютным эффектом и NNT. Ведущий сообщает, что статистическая значимость сама по себе не определяет клиническую важность, и следующий слайд проверяет применимость результата к конкретному пациенту.

\subsection*{Слайд 85. Применимость результата к конкретному пациенту}
\begin{onslide}
На рисунке показан фрагмент обсуждения. Контрольные вопросы: соответствует ли пациент критериям включения; применим ли эффект в другой популяции; каковы ограничения подгруппы МРТ?
\end{onslide}
На рисунке показан фрагмент обсуждения. Контрольные вопросы включают соответствие пациента критериям включения, применимость эффекта в другой популяции и ограничения подгруппы МРТ. Ведущий указывает на выделенные фразы «generalisability» и «early stage PD». Проверка применимости завершает чтение статьи. Следующие два слайда собраны в расширенный чек-лист критического чтения.

\subsection*{Слайд 86. Расширенный чек-лист критического чтения РКИ}
\begin{onslide}
На слайде приведены пятнадцать пунктов проверки статьи: дизайн и PICO, рандомизация и сокрытие распределения, ослепление, ITT-анализ и потери, таблица 1, первичная конечная точка, величина эффекта и доверительный интервал, $p$-значение и множественные сравнения, графики, таблицы регрессии, пропуски, нежелательные явления, подгруппы, ограничения и применимость результата.
\end{onslide}
На слайде приведены пятнадцать пунктов проверки статьи. Чек-лист охватывает дизайн и PICO, рандомизацию и сокрытие распределения, ослепление, ITT-анализ и потери, таблицу 1 и первичную конечную точку. Далее проверяются величина эффекта и доверительный интервал, $p$-значение и множественные сравнения, графики, таблицы регрессии, пропуски, нежелательные явления, подгруппы, ограничения и применимость результата. Ведущий сообщает, что чек-лист завершает блок чтения клинической статьи. Пятнадцать пунктов охватывают путь от дизайна до клинической интерпретации. Следующий слайд содержит ответы для статьи Johansson et al.

\subsection*{Слайд 87. Ответы чек-листа: дизайн и отчётность}
\begin{onslide}
На слайде приведены ответы на пятнадцать пунктов чек-листа. Исследование является двойным слепым РКИ при ранней стадии болезни Паркинсона. Аэробную нагрузку сравнивали с растяжкой как активным контролем. Распределение выполнено 1:1, участники и оценщики были ослеплены. Сведения о генерации последовательности и сокрытии распределения отсутствуют. Поток участников приведён с причинами выбытия. Таблица 1 описывает сопоставимость групп. В таблице 3 статьи $g$ Хеджеса равно 1,29 при 95\% доверительном интервале [0,65; 1,91] для контраста «аэробная нагрузка против растяжки» в кластере правой DLPFC; показатель стандартизует z-статистику, усреднённую по кластеру. Корреляция $\rho = 0{,}62$, $p = 0{,}003$ получена в подгруппе аэробной нагрузки ($n=21$) и носит поисковый характер. Часть анализов скорректирована на множественные сравнения, часть носит поисковый характер. Графики содержат подписи и единицы. Таблица 3 содержит оценки эффекта. Пропуски в подгруппе МРТ описаны в потоке участников. В учебном датасете шесть пропусков находятся в колонках «ca» и «thal», которые в демонстрациях не используются. В рассматриваемой публикации систематизированный перечень нежелательных явлений не приведён. Подгруппа МРТ формировалась без случайного отбора. Результаты применимы к ранней стадии болезни с оговоркой на состав подгруппы МРТ.
\end{onslide}
На слайде приведены ответы на пятнадцать пунктов чек-листа для статьи Johansson et al. Исследование является двойным слепым РКИ при ранней стадии болезни Паркинсона, а аэробную нагрузку сравнивали с растяжкой как активным контролем; распределение выполнено 1:1, участники и оценщики были ослеплены, но сведения о генерации последовательности и сокрытии распределения отсутствуют. Поток участников приведён с причинами выбытия, таблица 1 описывает сопоставимость групп, а величина эффекта $g$ Хеджеса равна 1,29 при 95\% доверительном интервале [0,65; 1,91]. Корреляция $\rho = 0{,}62$, $p = 0{,}003$, часть анализов скорректирована на множественные сравнения, а часть носит поисковый характер; графики содержат подписи и единицы. Таблица 3 содержит оценки эффекта, пропуски в подгруппе МРТ описаны в потоке участников, а систематизированный перечень нежелательных явлений отсутствует. Подгруппа МРТ формировалась без случайного отбора, а результаты применимы к ранней стадии болезни с оговоркой на состав подгруппы МРТ; ведущий указывает, что ответы не расширяют вывод за пределы дизайна исследования и что чтение статьи завершается переходом к разделу об ИИ.

\subsection*{Слайд 88. Ответы чек-листа: результаты и ограничения}
\begin{onslide}
На слайде приведены ответы на пункты 9--15 чек-листа: графики и подписи, таблица 3, пропуски в подгруппе МРТ, нежелательные явления, подгруппы, ограничения и применимость результата.
\end{onslide}
Ответы показывают, что интерпретация статьи требует разделения дизайна, результатов и ограничений. В рассматриваемой публикации систематизированный перечень нежелательных явлений не приведён. Применимость результата ограничена ранней стадией болезни и составом подгруппы МРТ. Следующий раздел переходит к вопросам применения ИИ.
\subsection*{Слайд 89. ИИ}
\begin{onslide}
Слайд-разделитель с названием раздела.
\end{onslide}
На слайде показан разделитель раздела «ИИ». Ведущий сообщает, что перед завершением фиксируются границы применения ИИ. Разделитель подготавливает переход от чтения статьи к ограничениям инструмента. Следующий слайд описывает риск недостоверной генерации.

\subsection*{Слайд 90. Риск недостоверной генерации}
\begin{onslide}
На слайде указано, что ИИ может уверенно описывать несуществующие функции и неверные методы. Ниже указано, что код или метод, предложенный ИИ, сначала проверяется на учебном наборе, затем сверяется с документацией.
\end{onslide}
На слайде указано, что ИИ может уверенно описывать несуществующие функции и неверные методы. Ниже указано, что код или метод, предложенный ИИ, сначала проверяется на учебном наборе, затем сверяется с документацией. Ведущий поясняет, что проверка на учебном наборе и сверка с документацией позволяют выявить неверный метод до его применения в анализе. Следующий слайд уточняет роль ИИ как вспомогательного инструмента.

\subsection*{Слайд 91. ИИ служит вспомогательным инструментом}
\begin{onslide}
На слайде указано, что ИИ помогает объяснить таблицу и подготовить шаблон анализа; интерпретацию результата выполняет исследователь. Блок «Возможности Jamovi» сообщает, что по состоянию на 2026 год официальный ИИ-помощник в классическом Jamovi отсутствует. Сторонний модуль askLLM доступен по адресу \url{https://scchen.com/project/askllm-jamovi/}; в заседании он не демонстрируется.
\end{onslide}
На слайде указано, что ИИ помогает объяснить таблицу и подготовить шаблон анализа, а интерпретацию результата выполняет исследователь. Блок «Возможности Jamovi» сообщает, что по состоянию на 2026 год официальный ИИ-помощник в классическом Jamovi отсутствует. Сторонний модуль askLLM доступен по адресу \url{https://scchen.com/project/askllm-jamovi/}, но в заседании он не демонстрируется. Ведущий сообщает, что ответственность за вывод остаётся за исследователем. Следующий раздел подводит итоги заседания.

\subsection*{Слайд 92. Итоги}
\begin{onslide}
Слайд-разделитель с названием раздела.
\end{onslide}
На слайде показан разделитель раздела «Итоги». Ведущий сообщает, что блок резюмирует приобретённые умения. Разделитель завершает содержательную часть заседания. Следующий слайд содержит три основных вывода.

\subsection*{Слайд 93. Три основных вывода заседания}
\begin{onslide}
Три вывода: Jamovi обеспечивает базовый анализ без программирования: импорт, описание, тесты и графики; принцип чтения $p$-значения и доверительного интервала сохраняется во всех таблицах; выбор критерия определяется вопросом, типом данных, дизайном и предпосылками модели. Ниже указано, что датасет и презентация находятся в репозитории заседания и вопросы задаются сейчас.
\end{onslide}
На слайде приведены три вывода. Jamovi обеспечивает базовый анализ без программирования: импорт, описание, тесты и графики. Принцип чтения $p$-значения и доверительного интервала сохраняется во всех таблицах. Выбор критерия определяется типом данных, типом групп и нормальностью распределения. Ниже указано, что датасет и презентация находятся в репозитории заседания и что вопросы задаются сейчас. Ведущий добавляет, что байесовский ответ дополняет частотный, а чтение статьи начинается с дизайна; следующий слайд содержит список литературы.

\subsection*{Слайд 94. Литература и источники}
\begin{onslide}
Приведены три колонки ссылок. Основные: Vu J., Harrington D. Introductory Statistics for the Life and Biomedical Sciences. OpenIntro, 2021; Петри А., Сэбин К. Наглядная медицинская статистика. ГЭОТАР-Медиа, 2021; Гланц С. Медико-биологическая статистика. Практика, 1999; документация Jamovi. Байесовский статистический вывод: Rouder J. et al. Bayesian $t$ tests. Psychon Bull Rev, 2009; Kruschke J. J Exp Psychol Gen, 2013; Wagenmakers E.-J. et al. 2018; Greenland S. et al. Eur J Epidemiol, 2016. Чтение статьи: Hopewell S. et al. CONSORT 2025; von Elm E. et al. STROBE 2007; Greenhalgh T. BMJ, 1997; Johansson M. et al. Ann Neurol, 2022; датасет UCI Heart Disease (1989).
\end{onslide}
На слайде приведены три колонки ссылок. Основные источники включают Vu J., Harrington D. Introductory Statistics for the Life and Biomedical Sciences, OpenIntro, 2021; Петри А., Сэбин К. Наглядная медицинская статистика, ГЭОТАР-Медиа, 2021; Гланц С. Медико-биологическая статистика, Практика, 1999 и документацию Jamovi. Раздел по байесовскому статистическому выводу содержит работы Rouder J. et al., 2009, Kruschke J., 2013, Wagenmakers E.-J. et al., 2018 и Greenland S. et al., 2016. Раздел по чтению статьи содержит CONSORT 2025 Hopewell S. et al., STROBE 2007 von Elm E. et al., Greenhalgh T., BMJ, 1997, Johansson M. et al., Ann Neurol, 2022 и датасет UCI Heart Disease (1989). Локальные копии находятся в каталоге \texttt{literature/}. Следующий слайд завершает заседание контактной информацией.

\subsection*{Слайд 95. Спасибо за внимание}
\begin{onslide}
На слайде приведены контактная информация ведущего и логотип вуза.
\end{onslide}
На слайде приведены контактная информация ведущего и логотип вуза. Ведущий сообщает, что вопросы задаются сейчас. Слайд завершает заседание и оставляет время для обсуждения. Материалы заседания доступны в репозитории, указанном на слайде 2.

\section*{Источники формул}

Формулы сверены с академическими и официальными источниками.

\begin{itemize}\footnotesize
\item Описательные показатели и преобразования: OpenIntro, \emph{Introductory Statistics for the Life and Biomedical Sciences}, 2021; \url{https://openintro.org/book/biostat}.
\item Асимметрия, эксцесс и проверка нормальности: Altman D. G., Bland J. M. Statistics Notes: the normal distribution. \emph{BMJ}, 1995;Ghasemi A., Zahediasl S.. Normality tests for statistical analysis. \url{https://pmc.ncbi.nlm.nih.gov/articles/PMC3693611/}.
\item $t$-критерии, Уэлч, Левен, доверительный интервал, мощность: Student. The probable error of a mean. \emph{Biometrika}, 1908; Welch B. L. The generalization of Student's ratio. \emph{Annals of Mathematical Statistics}, 1947.
\item Корреляция Пирсона, Спирмена и преобразование Фишера: Fisher R. A. Frequency distribution of the values of the correlation coefficient. \emph{Biometrika}, 1915.
\item Хи-квадрат и $V$ Крамера: Cramér H. \emph{Mathematical Methods of Statistics}, 1946.
\item Байесовский фактор и формула Байеса: Kass R. E., Raftery A. E. Bayes Factors. \emph{JASA}, 1995; Wagenmakers E.-J. et al. Bayesian inference for psychology, 2018.
\item Размер эффекта $d$ Коэна и $g$ Хеджеса: Cohen J. \emph{Statistical Power Analysis for the Behavioral Sciences}, 1988; Hedges L. V. Distribution theory for Glass's estimator of effect size. \emph{JESA}, 1981.
\item Логистическая регрессия, OR и NNT: Cochrane Handbook, chapter 15; \url{https://www.cochrane.org/authors/handbooks-and-manuals/handbook/current/chapter-15}.
\item ROC-кривая: Hanley J. A., McNeil B. J. The meaning and use of the area under a ROC curve. \emph{Radiology}, 1982.
\item Оценка Каплана--Мейера: Kaplan E. L., Meier P. Nonparametric estimation from incomplete observations. \emph{JASA}, 1958.
\end{itemize}

\appendix
\section*{Терминологический указатель}
\begin{description}
\item[Байесовский статистический вывод] Метод вывода, в котором неизвестные параметры
и гипотезы получают распределения вероятностей. Корректное название метода.
\item[Байесовский фактор] Отношение правдоподобий данных при двух гипотезах;
обозначается $\mathrm{BF}_{10}$ или $\mathrm{BF}_{01}$.
\item[Априорное распределение] Распределение, описывающее представления до наблюдения данных.
\item[Апостериорное распределение] Распределение, обновлённое с учётом данных.
\item[Апостериорный интервал] Интервал, содержащий указанную долю апостериорной вероятности;
английский термин credible interval, сокращение CrI.
\item[ROPE] Область практической эквивалентности, диапазон пренебрежимо малых значений эффекта.
\item[ITT-анализ] Анализ всех рандомизированных пациентов в их исходных группах.
\item[Сокрытие распределения] Процедура, исключающая раскрытие группы до рандомизации;
английский термин allocation concealment.
\item[Ошибка первого рода для семейства гипотез] Английский термин family-wise error rate,
сокращение FWER.
\item[g Хеджеса] Стандартизованная величина эффекта; английский термин Hedges' g.
\item[V Крамера] Мера силы связи между категориальными признаками.
\item[Оценка Каплана--Мейера] Непараметрическая оценка функции выживаемости.
\item[Форест-плот] График оценок эффектов исследований с доверительными интервалами.
\item[PPV] Прогностическая ценность положительного результата.
\end{description}

\section*{Подготовка ведущего}
\begin{itemize}
\item Проверить импорт файла \texttt{heart-disease.csv} и типы переменных на своей машине.
\item Проверить читаемость QR-кодов с задних рядов.
\item Проверить, что PDF презентации открывается и содержит 93 слайда.
\item Проверить наличие локальных копий литературы в каталоге \texttt{literature/}.
\end{itemize}

\end{document}
